% Opcje klasy 'iithesis' opisane sa w komentarzach w pliku klasy. Za ich pomoca
% ustawia sie przede wszystkim jezyk oraz rodzaj (lic/inz/mgr) pracy.


\documentclass[shortabstract]{iithesis}

\usepackage[utf8]{inputenc}
\usepackage[T1]{fontenc}
\usepackage{lmodern}
\usepackage{longtable}
\usepackage{amsmath,amssymb}
\usepackage{amsthm}
\usepackage{caption}
\usepackage{bussproofs}
\usepackage{algorithm}
\usepackage{algpseudocode}
\usepackage{listings}
\usepackage{float}
\raggedbottom
\usepackage{enumitem}
\lstset{
  basicstyle=\ttfamily\small,
  columns=fullflexible,
  keepspaces=true,
  literate=
    {∧}{{$\wedge$}}1
    {∨}{{$\vee$}}1
    {¬}{{$\neg$}}1
    {→}{{$\to$}}1
    {↔}{{$\leftrightarrow$}}1
    {⊤}{{$\top$}}1
    {⊥}{{$\bot$}}1
    {⊢}{{$\vdash$}}1
}
\theoremstyle{definition}
\newtheorem{definition}{Definicja}

\theoremstyle{remark}
\newtheorem{observation}[definition]{Obserwacja}

\theoremstyle{plain}
\newtheorem{theorem}[definition]{Twierdzenie}
\newtheorem{lemma}[definition]{Lemat}
\newtheorem{proposition}[definition]{Wniosek}
%%%%% DANE DO STRONY TYTUŁOWEJ
% Niezaleznie od jezyka pracy wybranego w opcjach klasy, tytul i streszczenie
% pracy nalezy podac zarowno w jezyku polskim, jak i angielskim.
% Pamietaj o madrym (zgodnym z logicznym rozbiorem zdania oraz estetyka) recznym
% zlamaniu wierszy w temacie pracy, zwlaszcza tego w jezyku pracy. Uzyj do tego
% polecenia \fmlinebreak.
\polishtitle    {Wymagający złamania wierszy\fmlinebreak tytuł pracy w~języku polskim}
\englishtitle   {English title}
\polishabstract {\ldots}
\englishabstract{\ldots}
% w pracach wielu autorow nazwiska mozna oddzielic poleceniem \and
\author         {Maksymilian Debeściak}
% w przypadku kilku promotorow, lub koniecznosci podania ich afiliacji, linie
% w ponizszym poleceniu mozna zlamac poleceniem \fmlinebreak
\advisor        {dr Jan Kowalski}
%\date          {}                     % Data zlozenia pracy
% Dane do oswiadczenia o autorskim wykonaniu
%\transcriptnum {}                     % Numer indeksu
%\advisorgen    {dr. Jana Kowalskiego} % Nazwisko promotora w dopelniaczu
%%%%%

%%%%% WLASNE DODATKOWE PAKIETY
%
%\usepackage{graphicx,listings,amsmath,amssymb,amsthm,amsfonts,tikz}
%
%%%%% WŁASNE DEFINICJE I POLECENIA
%
%\theoremstyle{definition} \newtheorem{definition}{Definition}[chapter]
%\theoremstyle{remark} \newtheorem{remark}[definition]{Observation}
%\theoremstyle{plain} \newtheorem{theorem}[definition]{Theorem}
%\theoremstyle{plain} \newtheorem{lemma}[definition]{Lemma}
%\renewcommand \qedsymbol {\ensuremath{\square}}
% ...
%%%%%
%%%%% POCZĄTEK ZASADNICZEGO TEKSTU PRACY

\begin{document}




\chapter{Wprowadzenie}
Dedukcja naturalna jest metodą dowodzenia twierdzń logicznych przez wychodzenie z pewnych założeń i modyfikowanie ich za pomocą
ściśle zdefiniowanych reguł wnioskowania. 
\section{Dedukcja naturalna}
\begin{definition}
\textbf{Sekwentem} nazywamy wyrażenie postaci

\[
\Gamma \vdash \varphi
\]
gdzie $\Gamma$ jest skończonym zbiorem formuł (zapisywanych po przecinku i bez nawiasów klamrowych) zwanym \emph{kontekstem},
a $\varphi$ jest formułą zwaną \emph{tezą sekwentu}. Sekwent $\Gamma \vdash \varphi$ interpretujemy jako stwierdzenie, że jeśli założymy, że wszystkie formuły z $\Gamma$ są prawdziwe, 
da się dowieść $\varphi$. W rachunku zdań w logice klasycznej, który będzie jedyną logiką w której się będziemy poruszać, zdanie prawdziwe we wszystkich wartościowaniach jest dowodliwe, tak więc $\Gamma \vdash \phi$ jest równoważne stwierdzeniu jeżeli wszystkie formuły kontekstu $\Gamma$ są prawdziwe to $\varphi$ jest prawdziwe.
\end{definition}


\begin{definition}
    \textbf{Regułą wnioskowania} nazywamy schemat postaci
\[
\frac{S_1 \qquad S_2 \qquad \dots \qquad S_n}{S}(R)
\]
gdzie $S_1,\dots,S_n,S$ są sekwentami. Sekwenty $S_1,\dots,S_n$ nazywamy \emph{przesłankami reguły}, natomiast sekwent $S$ nazywamy \emph{konkluzją reguły}. $R$ to z kolei \emph{indeks reguły}, który służy do jej identyfikacji. 
Jeżeli przesłanki reguły są prawdziwe to z definicji jej konkluzja jest również prawdziwa.
\end{definition}

Dowody w dedukcji naturalnej można utożsamiać z drzewami, których węzłami są sekwenty,
 a krawędzie są indeksowane indeksami reguł wnioskowania.

W nieniejszej pracy skupimy się na dowodzeniu twierdzeń w logice klasycznej z następującymi spójnikiami logicznymi: koniunkcją ($\land$), alternatywą ($\lor$), negacją ($\neg$) oraz implikacją ($\to$).
regułami wnioskowania:

\begin{center}
\[
\begin{array}{c c}
\displaystyle
\overline{\ \varphi \vdash \varphi\ }\;(\mathrm{Ass})
&
\displaystyle
\frac{\Gamma \vdash \varphi}{\Gamma,\psi \vdash \varphi}\;(\mathrm{W})
\\[2.2ex]
\displaystyle
\frac{\Gamma,\varphi \vdash \psi}{\Gamma \vdash \varphi \to \psi}\;(\to I)
&
\displaystyle
\frac{\Gamma_1 \vdash \varphi \to \psi \qquad \Gamma_2 \vdash \varphi}
     {\Gamma_1,\Gamma_2 \vdash \psi}\;(\to E)
\\[2.2ex]
\displaystyle
\frac{\Gamma,\varphi \vdash \bot}{\Gamma \vdash \neg\varphi}\;(\neg I)
&
\displaystyle
\frac{\Gamma_1 \vdash \neg\varphi \qquad \Gamma_2 \vdash \varphi}
     {\Gamma_1,\Gamma_2 \vdash \bot}\;(\neg E)
\\[2.2ex]
\displaystyle
\frac{\Gamma_1 \vdash \varphi \qquad \Gamma_2 \vdash \psi}
     {\Gamma_1,\Gamma_2 \vdash \varphi \wedge \psi}\;(\wedge I)
&
\displaystyle
\frac{\Gamma \vdash \varphi \wedge \psi}{\Gamma \vdash \varphi}\;(\wedge E_1)
\\[2.2ex]
\displaystyle
\frac{\Gamma \vdash \varphi}{\Gamma \vdash \varphi \vee \psi}\;(\vee I_1)
&
\displaystyle
\frac{\Gamma \vdash \varphi \wedge \psi}{\Gamma \vdash \psi}\;(\wedge E_2)
\\[2.2ex]
\displaystyle
\frac{\Gamma \vdash \psi}{\Gamma \vdash \varphi \vee \psi}\;(\vee I_2)
&
\displaystyle
\frac{\Gamma_1 \vdash \varphi \vee \psi \qquad \Gamma_2,\varphi \vdash \rho \qquad \Gamma_3,\psi \vdash \rho}
     {\Gamma_1,\Gamma_2,\Gamma_3 \vdash \rho}\;(\vee E)
\\[2.2ex]
\displaystyle
\overline{\ \vdash \top\ }\;(\top I)
&
\displaystyle
\frac{\Gamma \vdash \bot}{\Gamma \vdash \varphi}\;(\bot E)
\\[2.2ex]
\displaystyle
\frac{\Gamma_1,\varphi \vdash \psi \qquad \Gamma_2,\psi \vdash \varphi}
     {\Gamma_1,\Gamma_2 \vdash \varphi \leftrightarrow \psi}\;(\leftrightarrow I)
&
\displaystyle
\frac{\Gamma_1 \vdash \varphi \leftrightarrow \psi \qquad \Gamma_2 \vdash \varphi}
     {\Gamma_1,\Gamma_2 \vdash \psi}\;(\leftrightarrow E_1)
\\[2.2ex]
\displaystyle
\frac{\Gamma_1 \vdash \varphi \leftrightarrow \psi \qquad \Gamma_2 \vdash \psi}
     {\Gamma_1,\Gamma_2 \vdash \varphi}\;(\leftrightarrow E_2)
&
\displaystyle
\frac{\Gamma,\neg\varphi \vdash \bot}{\Gamma \vdash \varphi}\;(\mathrm{RAA})
\\[2.2ex]
\displaystyle
\frac{\Gamma \vdash \neg\neg\varphi}{\Gamma \vdash \varphi}\;(\neg\neg)
&
\displaystyle
\overline{\ \vdash \varphi \vee \neg\varphi\ }\;(\mathrm{TND})

\end{array}
\]
\label{tab:rules}
\captionof{table}{Reguły wnioskowania rozważanego systemu}



\end{center}
Oto przykładowy dowód w dedukcji naturalnej w notacji drzewiastej:


\begin{prooftree}

% Ass: a ∨ b ⊢ a ∨ b
\AxiomC{}
\RightLabel{$(\mathrm{Ass})$}
\UnaryInfC{$a \vee b \vdash a \vee b$}

% Ass: a ⊢ a
\AxiomC{}
\RightLabel{$(\mathrm{Ass})$}
\UnaryInfC{$a \vdash a$}
\RightLabel{$(\vee I_2)$}
\UnaryInfC{$a \vdash b \vee a$}

% Ass: b ⊢ b
\AxiomC{}
\RightLabel{$(\mathrm{Ass})$}
\UnaryInfC{$b \vdash b$}
\RightLabel{$(\vee I_1)$}
\UnaryInfC{$b \vdash b \vee a$}

% ∨E
\RightLabel{$(\vee E)$}
\TrinaryInfC{$a \vee b \vdash b \vee a$}

% →I
\RightLabel{$(\to I)$}
\UnaryInfC{$\vdash (a \vee b) \to (b \vee a)$}

\end{prooftree}


Dużą zaletą dedukcji naturalnej jest to, że niejako symuluje ona sposób, w jaki ludzie przeprowadzają wnioskowanie, zarówno 
w matematyce, jak i w codziennym życiu. Dowody mają swoją drzweiastą strukturę, którą można łatwo analizować.
Inną zaletą tej metody jest jej uniwersalność. Dodając odpowiednie reguły wnioskowania
możemy otrzymać dedukcję naturalną dla różnych systemów logicznych np logiki pierwszego i drugiego rzędu, czy logiki modalnej.
Jeszcze jedną zaletą jest to, że w przypadku większości użytecznych dowodów jesteśmy w stanie uniknąć wykładniczo dużej złożoności 
wielu innych systemów dowodzenia sprowadzających się do rozpatrzenia wszystkich możliwych przypadków.
Z drugiej strony dedukcja naturalna jest trudna w zaitomatyzowaniu. Zwykle dowód przeprowadzamy z góry na dół, zaczynając od założeń i
kończąc na tezie. Jednak nie zawsze jasne jest, którą regułę wnioskowania zastosować i co gorsza wyniki zastosowania niektórych reguł mogą zawierać
podformuły których nie ma w przesłankach, lub jej wybór nie jest jednoznaczny. Np używając reguły 

\[
\displaystyle
\frac{\Gamma \vdash \varphi}{\Gamma \vdash \varphi \vee \psi}\;(\vee I_1)
\]

w przesłance nie mamy informacji o formule $\psi$, która pojawi się w konkluzji. 
Jednak podobieństwo dedukcji naturalnej do ludzkiego sposobu wnioskowania sprawia, oraz pewien dający się zauważyć porządek
w zapisie dowodów sprawia, że gdzie klasyczna implementacja dedukcji naturalnej jest trudna, wydaje się ona dobrze nadawać
do implementacji z użyciem metod sztucznej inteligencji, w tym modeli językowych o arhirtekturze transformerów.

W niniejszej pracy przedstawimy implementację automatycznego systemu dowodzenia twierdzeń w logice klasycznej przy użyciu odpowiednio 
wytrenowanego modelu nano-GPT oraz przedstawimy wyniki jego działania i jak dobrze radzi sobie z przeprowadzaniem dowodów. Konstrukcja
systemu będzie opisana krok po kroku, zaczynając od zaprogrogramowania logicznej struktury systemu, tłumaczenia zapisu dowodu w formie 
tekstowej na strukturę logiczną i sprawdzanie poprawności dowodu, generacja korpusu treningowego, trenowanie modelu oraz 
kolejne etapy konstrukcj samego systemu dowodzenia(wybór możliwych linii, równoległość prowadzenia kilku ścieżek rozumowania, itp).


\section{Zapis tekstowy dowodów}

Notacja Ficha jest sposobem zapisu dowodów dedukcji naturalnej w formie tekstowej. Jest on wygodny, ponieważ metoda zapisywania drzew dowodu potrafi zajmować bardzo dużo miejsca
i trudno użyć do niej modeli językowych. 

Z punktu widzenia implementacji komputerowej klasyczna notacja Fitcha narzuca
ograniczenia wynikające ze struktury zagnieżdżonych bloków, które określają,
do których założeń można się odwoływać w danym miejscu dowodu.
W konsekwencji dostęp do wcześniejszych kroków dowodu zależy od ich położenia
w tej strukturze.

W stosowanej w niniejszej pracy notacji dowód koduje uporządkowaną
lista sekwentów w,której każda linia koduje jeden sekwent, a zależności między kolejnymi krokami są wyrażane poprzez odwołania
do indeksów wcześniejszych linii oraz zastosowane reguły wnioskowania.
Nie występują tu ograniczenia wynikające ze zagnieżdżonych bloków, a dowolny wcześniejszy
sekwent może zostać użyty jako przesłanka, o ile spełnia warunki danej reguły.

Takie podejście pozwala traktować konstrukcję dowodu jako proces składania sekwentów
na podstawie wcześniej uzyskanych wyników, bez konieczności odtwarzania struktury
zagnieżdżeń charakterystycznej dla klasycznej notacji Fitcha.

Co więcej, w stosowanej notacji każdej linii zapisywana jest jawnie teza sekwentu, który dana linia koduje.
Jest to użyteczne podczas sprawdzania poprawności dowodu.

Podejście to wprowadza jednak pewną niedogodność: nie istnieje jedna reguła, pozwalająca wyciągnąć formułę należącą do kontekstu
sekwentu w jednym kroku. Ponieważ jest to dość częsta operacja wprowadzamy dwie pomocnicze reguły, których celem nie jest wprowadzanie
 nowych własności logicznych systemu, lecz stanowią wygodny mechanizm umożliwiający prostsze manipulowanie kontekstem.


\[
\displaystyle
\frac{\varphi , \Gamma \vdash \psi}{\Gamma, \varphi \vdash \varphi}\;(\mathrm{FromContext})
\quad
\displaystyle
\frac{\Gamma \vdash \psi}{\Gamma , \varphi \vdash \varphi}\;(\mathrm{FromWeakenContext})
\]


Poniżej przedstawiam drzewa dowodu poprawności tych reguł:




\begin{prooftree}
\AxiomC{$\Gamma, \varphi \vdash \psi$}
\AxiomC{}
\RightLabel{$(\mathrm{Ass})$}
\UnaryInfC{$\varphi \vdash \varphi$}
\RightLabel{$(\land I)$}
\BinaryInfC{$\Gamma, \varphi \vdash \psi \land \varphi$}
\RightLabel{$(\land E_2)$}
\UnaryInfC{$\Gamma, \varphi \vdash \varphi$}
\end{prooftree}
\begin{center}
\textbf{Dowód reguły $(\mathrm{FromContext})$}
\end{center}



\begin{prooftree}
\AxiomC{$\Gamma \vdash \psi$}
\AxiomC{}
\RightLabel{$(\mathrm{Ass})$}
\UnaryInfC{$\varphi \vdash \varphi$}
\RightLabel{$(\land I)$}
\BinaryInfC{$\Gamma, \varphi \vdash \psi \land \varphi$}
\RightLabel{$(\land E_2)$}
\UnaryInfC{$\Gamma, \varphi \vdash \varphi$}
\end{prooftree}
\begin{center}
\textbf{Dowód reguły $(\mathrm{FromWeakenContext})$}
\end{center}

Aby sformalizować podejście, w którym każda linia dowodu koduje jeden sekwent, wprowadzamy oznaczenia pozwalające jednoznacznie
interpretować poszczególne linie dowodu oraz odwołania między nimi:


$$\texttt{Fitch(i)} :=\text{sekwent kodowany przez linię o indeksie } \texttt{i}$$ 
oraz oznaczenie
$$\texttt{s.conclusion} := \text{teza sekwentu } \texttt{s}$$
$$\texttt{s.assumptions} := \text{kontekst sekwentu } \texttt{s}$$

Poniżej przedstawimy interpretację poszczególnych reguł wnioskowania w notacji Ficha, a także ich odpowiedniki w notacji drzewiastej:

\begin{center}
\small
\begin{longtable}{|p{0.42\textwidth}|p{0.48\textwidth}|}
\hline
\textbf{Notacja Fitcha} & \textbf{Notacja drzewiasta} \\
\hline

\texttt{i. $\phi$    assumption}
%DODAC REGULE CONTEXT
&
$$\displaystyle
\overline{\ \texttt{$\phi\vdash\phi$}\ }\;(\mathrm{Ass})
$$
\\
\hline
\texttt{j. $\psi \quad \dots$
     \newline
     \dots
     \newline
     k. $\phi \quad \dots$
     \newline
     \dots
     \newline
     i. $\phi$ \quad   weakening, j, k} %TO NAPRAWIĆ!!!!!!!!!!!!!!!!!!!!!!!!!!!!!!!!!!!!!!!!!
&
$$\displaystyle
\frac{\texttt{Fitch(k)}}{\texttt{Fitch(i)}}\;(W)
$$
gdzie 
\texttt{Fitch(i).assumptions = \newline Fitch(k).assumptions, \newline Fitch(j).conclusion } 

\\
\hline
\texttt{j. $\phi \quad  \dots$
     \newline
     $\dots$
     \newline
     k. $\psi \quad \dots$ 
     \newline
     $\dots$
     \newline
     i. $\phi \to \psi \quad$ →-introduction, j–k
     }





&
$$\displaystyle
\frac{\texttt{Fitch(k)}}{\texttt{Fitch(i)}}\;(\to I)
$$
\\
\hline

\texttt{j. $\phi \quad  \dots$
     \newline
     $\dots$
     \newline
     k. $\bot \quad \dots$ 
     \newline
     $\dots$
     \newline
     i. $\neg \phi \quad$ ¬-introduction, j–k
     }





&
$$\displaystyle
\frac{\texttt{Fitch(k)}}{\texttt{Fitch(i)}}\;(\neg I)
$$
\\









\hline

\texttt{j. $\phi \quad  \dots$
     \newline
     $\dots$
     \newline
     k.  $\phi \to \psi \quad  \dots$
     \newline
     $\dots$
     \newline
     i. $\psi \qquad$ →-elimination, k, j
     }





&
$$\displaystyle
\frac{\texttt{ Fitch(k) \quad Fitch(j)}}{\texttt{Fitch(i)}}\;(\to E)
$$
\\

\hline

\texttt{j. $ \neg \phi \quad  \dots$
     \newline
     $\dots$
     \newline
     k.  $\phi  \quad  \dots$
     \newline
     $\dots$
     \newline
     i. $\bot \quad$ ¬-elimination, j, k
     }





&
$$\displaystyle
\frac{\texttt{ Fitch(j) \qquad Fitch(k)}}{\texttt{Fitch(i)}}\;(\neg E)
$$
\\

\hline





\texttt{j. $  \phi \quad  \dots$
     \newline
     $\dots$
     \newline
     k.  $\psi  \quad  \dots$
     \newline
     $\dots$
     \newline
     i. $\phi \land \psi \quad$ $\wedge$-introduction, j, k}





&
$$\displaystyle
\frac{\texttt{ Fitch(j) \quad Fitch(k)}}{\texttt{Fitch(i)}}\;(\wedge I)
$$






\\

\hline



\texttt{j. $  \phi \quad  \dots$
     \newline
     $\dots$
     \newline
     i. $\phi \lor \psi \quad$ $\lor$-introduction1, j}
&
$$\displaystyle
\frac{\texttt{Fitch(j)}}{\texttt{Fitch(i)}}\;(\lor I_1)
$$
\\

\hline


\texttt{j. $  \psi \quad  \dots$
     \newline
     $\dots$
     \newline
     i. $\phi \lor \psi \quad$ $\lor$-introduction2, j}
&
$$\displaystyle
\frac{\texttt{Fitch(j)}}{\texttt{Fitch(i)}}\;(\lor I_2)
$$
\\

\hline








\texttt{i. $\top \quad \top$-introduction}
&
$$\displaystyle
\overline{\ \vdash \top\ }\;(\top I)$$
\\

\hline


\texttt{j. $  \phi \land \psi  \quad  \dots$
     \newline
     $\dots$
     \newline
     i. $\phi \quad$ $\land$-elimination1, j}
&
$$\displaystyle
\frac{\texttt{Fitch(j)}}{\texttt{Fitch(i)}}\;(\land E_1)
$$
\\

\hline



\texttt{j. $  \phi \land \psi  \quad  \dots$
     \newline
     $\dots$
     \newline
     i. $\psi \quad$ $\land$-elimination2, j}
&
$$\displaystyle
\frac{\texttt{Fitch(j)}}{\texttt{Fitch(i)}}\;(\land E_2)
$$
\\

\hline










\texttt{j. $  \phi \lor \psi  \quad  \dots$
     \newline
     $\dots$
     \newline
     k.  $\rho  \quad  \dots$
     \newline
     $\dots$
     \newline
     l. $\rho  \quad  \dots$
     \newline
     $\dots$
     \newline     i. $\rho \quad$ $\lor$-elimination, j, k, l}
&
$$
\displaystyle
\frac{\texttt{Fitch(j)} \qquad \texttt{Fitch(k)} \qquad \texttt{Fitch(l)}}
     {\texttt{Fitch(i)}}\;(\vee E)
$$
     \\



\hline

\texttt{j. $  \bot \quad  \dots$
     \newline
     $\dots$
     \newline     i. $\phi \quad$ $\bot$-elimination, j}
&
$$
\displaystyle
\frac{\texttt{Fitch(j)} }
     {\texttt{Fitch(i)}}\;(\bot E)
$$
\\
\hline






\texttt{j. $\psi \quad  \dots$
     \newline
     $\dots$
     \newline
     k. $ \phi \quad  \dots$
     \newline
     $\dots$
     \newline
     i. $\phi \leftrightarrow \psi \quad $ $\leftrightarrow$-introduction, j, k}

&

$$\displaystyle
\frac{\texttt{Fitch(j)\qquad Fitch(k)}}{\texttt{Fitch(i)}}\;(\leftrightarrow I)
$$

\\
\hline


\texttt{j. $\phi \leftrightarrow \psi \quad  \dots$
     \newline
     $\dots$
     \newline
     k. $ \phi \quad  \dots$
     \newline
     $\dots$
     \newline
     i. $\psi  \quad $ $\leftrightarrow$-elimination1, j, k}

&

$$\displaystyle
\frac{\texttt{Fitch(j)\qquad Fitch(k)}}{\texttt{Fitch(i)}}\;(\leftrightarrow E_1)
$$

\\
\hline






\texttt{j. $\phi \leftrightarrow \psi \quad  \dots$
     \newline
     $\dots$
     \newline
     k. $ \psi \quad  \dots$
     \newline
     $\dots$
     \newline
     i. $\phi  \quad $ $\leftrightarrow$-elimination2, j, k}

&

$$\displaystyle
\frac{\texttt{Fitch(j)\qquad Fitch(k)}}{\texttt{Fitch(i)}}\;(\leftrightarrow E_2)
$$

\\
\hline



\texttt{j. $\neg \phi \quad  \dots$
     \newline
     $\dots$
     \newline
     k. $ \bot \quad  \dots$
     \newline
     $\dots$
     \newline
     i. $\phi  \quad $ RAA, j–k}

&

$$\displaystyle
\frac{\texttt{Fitch(k)}}{\texttt{Fitch(i)}}\;{(\mathrm{RAA})}
$$

\\
\hline



\texttt{j. $\neg \neg \phi \quad  \dots$
     \newline
     $\dots$
     \newline
     i. $\phi  \quad $ $\neg\neg$-elimination, j}

&

$$\displaystyle
\frac{\texttt{Fitch(j)}}{\texttt{Fitch(i)}}\;{(\neg\neg)}
$$

\\
\hline




\texttt{i. $\phi \lor \neg \phi \quad $TND}
&
$$\displaystyle
\overline{\ \vdash \phi \lor \neg \phi\ }\;{(\mathrm{TND})}$$
\\

\hline


\texttt{j. $\psi \quad \dots$
     \newline
     $\dots$
     \newline
     i. $\phi $ \quad context, j}
&
$$\displaystyle
\frac{\texttt{Fitch(j)}}{\texttt{Fitch(i)}}\;{(\mathrm{FromContext})}
$$
\\

\hline
\texttt{j. $\psi \quad \dots$
     \newline
     $\dots$
     \newline
     i. $\phi $ \quad contextWeak, j}
&
$$\displaystyle
\frac{\texttt{Fitch(j)}}{\texttt{Fitch(i)}}\;{(\mathrm{FromWeakenContext})}
$$
\\

\hline




\end{longtable}
\end{center}

Warto zauważyć, że aby dowody były poprawne, muszą być spełnione odpowiednie warunki
dotyczące sekwentów, do których się odwołujemy, wynikające z postaci reguł wnioskowania.
Warunki te nie są jednak podane explicite, lecz wynikają z możliwości dopasowania
struktury sekwentów do schematów reguł.



Na przykład dla zapisu:
\noindent
\begin{flushleft}
\texttt{j. $\phi \lor \psi$ \quad ... \\
... \\
k. $\rho$ \quad ... \\
... \\
l. $\rho$ \quad ... \\
... \\
i. $\rho$ \quad $\lor$-elimination, j, k, l
}
\end{flushleft}
który odpowiada dopasowaniu schematu
$$
\frac{\texttt{Fitch(j)} \qquad \texttt{Fitch(k)} \qquad \texttt{Fitch(l)}}
     {\texttt{Fitch(i)}}\;(\vee E)
$$
do schematu reguły 
$$
\frac{\Gamma_1 \vdash \varphi \vee \psi \qquad \Gamma_2,\varphi \vdash \rho \qquad \Gamma_3,\psi \vdash \rho}
     {\Gamma_1,\Gamma_2,\Gamma_3 \vdash \rho}\;(\vee E)
$$
poprawność zastosowania reguły $\lor E$ wymaga, aby teza sekwentu \texttt{Fitch(j)} była alternatywą, żeby
lewa strona tej alternatywy znajdowała się w kontekście sekwentu \texttt{Fitch(k)}, a prawa 
w kontekście sekwentu \texttt{Fitch(l)}, a także żeby tezy sekwentów \texttt{Fitch(k)} oraz \texttt{Fitch(l)} były takie same.  



Podobne warunki wynikają dla wszystkich reguł wnioskowania i sprowadzają się
do sprawdzenia, czy struktura użytych sekwentów odpowiada instancji danej reguły.

W związku z tym konieczne jest wprowadzenie formalizmu, który pozwala jednoznacznie
odtworzyć sekwent odpowiadający każdej linii dowodu, sprawdzić poprawność
ich dopasowania do reguł wnioskowania, a także zweryfikować, czy formuła zapisana
w danej linii odpowiada tezie kodowanego przez nią sekwentu.


\chapter{Reprezentacja formuł, sekwentów,reguł wnioskowania i dowodów}
\chaptermark{Reprezentacja systemu}

Pierwszym etapem konstrukcji systemu było zdefiniowanie reprezentacji obiektów
logicznych, na bazie których wykonywane będą dalsze operacje: parsowanie dowodów,
sprawdzanie poprawności zastosowań reguł, generowanie danych treningowych oraz
weryfikacja kroków proponowanych przez model językowy. Reprezentacja ta musiała
spełniać dwa podstawowe wymagania. Po pierwsze, powinna wiernie odpowiadać
formalizmowi dedukcji naturalnej opisanemu w poprzednim rozdziale. Po drugie,
powinna umożliwiać łatwe przekształcanie obiektów logicznych na postać tekstową,
ponieważ właśnie na takich tekstowych reprezentacjach uczony jest model językowy.

\section{Reprezentacja formuł i sekwentów}

Podstawowymi obiektami składniowymi są zmienne, relacje oraz formuły. W obecnej
wersji systemu wykorzystywane są jedynie zmienne zdaniowe, jednak reprezentacja
została zaprojektowana w sposób nieco ogólniejszy. Atom traktowany jest jako
zastosowanie relacji do listy argumentów. W przypadku rachunku zdań wystarcza
relacja jednoargumentowa, dlatego zapis taki jak \texttt{x} może być interpretowany
jako uproszczona postać atomu. Takie rozwiązanie nie jest konieczne dla samego
rachunku zdań, ale ułatwia ewentualne rozszerzenie systemu o bardziej ogólną
składnię, na przykład o predykaty wieloargumentowe.

Formuły złożone reprezentowane są rekurencyjnie. Formuła atomowa jest przypadkiem
bazowym, natomiast negacja, koniunkcja, alternatywa, implikacja i równoważność
przechowują swoje bezpośrednie podformuły. Dzięki temu struktura danych odpowiada
drzewiastej strukturze składniowej formuły. Taka reprezentacja pozwala zarówno
łatwo analizować postać formuły, jak i wykonywać operacje potrzebne przy stosowaniu
reguł wnioskowania, na przykład odczytanie lewej i prawej strony implikacji czy koniunkcji.

Kontekst dowodzenia reprezentuje zbiór formuł przyjętych jako założenia danego
sekwentu. W implementacji został on podzielony na dwie części: kontekst bazowy oraz
kontekst dodatkowy. Kontekst bazowy zawiera formuły traktowane jako założenia
obowiązujące w całym rozważanym fragmencie dowodu. Kontekst dodatkowy zawiera
natomiast formuły wprowadzane lokalnie przez reguły wnioskowania, na przykład
przy wprowadzaniu implikacji, negacji albo przy rozumowaniu nie wprost.

Podział ten nie zmienia znaczenia logicznego sekwentu, ale ma znaczenie techniczne
i algorytmiczne. Pozwala odróżnić założenia, które powinny być stale dostępne
podczas generowania poddowodu, od założeń lokalnych, które mogą być wprowadzane
i usuwane przez konkretne reguły. Jest to szczególnie istotne później, gdy system
będzie rozbijał jeden problem dowodowy na podproblemy oraz generował dowody
sekwentów z niepustym kontekstem.

Zdefiniowany wcześniej sekwent jest w systemie reprezentowany jako para złożona
z kontekstu oraz tezy. W implementacji obiekt ten nosi nazwę \texttt{LittleProblem},
ponieważ każdy sekwent można traktować jako mały problem dowodowy: należy
wyprowadzić określoną tezę przy założeniu określonego kontekstu. Taka nazwa jest
szczególnie naturalna w dalszej części pracy, gdzie większe problemy dowodowe będą
rozbijane na mniejsze podproblemy.

Reprezentacja sekwentu jest bezpośrednio wykorzystywana przez wszystkie dalsze
komponenty programu: reguły wnioskowania przekształcają listy sekwentów-przesłanek
w sekwenty-konkluzje, parser odtwarza sekwenty zakodowane w kolejnych liniach
dowodu, a weryfikator sprawdza, czy dana linia rzeczywiście odpowiada sekwentowi
otrzymanemu przez zastosowanie wskazanej reguły.
\section{Reprezentacja reguł wnioskowania}

Reguły wnioskowania zostały zaimplementowane jako osobne obiekty odpowiadające
regułom systemu przedstawionym w Tabeli~\ref{tab:rules}, a także pomocniczym
regułom \(\mathrm{FromContext}\) oraz \(\mathrm{FromWeakenContext}\). Na poziomie
programu każda reguła działa przede wszystkim jako funkcja częściowa: dla poprawnej
listy sekwentów-przesłanek zwraca sekwent będący konkluzją zastosowania tej reguły,
natomiast dla danych, które nie pasują do jej schematu, zgłasza błąd.

Można to opisać następującym schematem:


\begin{algorithm}[H]
\caption{Zastosowanie reguły wnioskowania}
\begin{algorithmic}[1]
\Procedure{ZastosujRegułę}{$R, P, A$}
    \If{$P$ nie spełnia warunków wejściowych reguły $R$}
        \State zgłoś błąd
    \EndIf
    \State $S \gets$ konkluzja reguły $R$ dla przesłanek $P$ i argumentów $A$
    \State \Return $S$
\EndProcedure
\end{algorithmic}
\end{algorithm}

Warunki wejściowe obejmują zgodność liczby przesłanek, zgodności ich typów (każda przesłanka musi być sekwentem), oraz ich zgodności ze schematem danej reguły. 
Argument \(R\) oznacza regułę wnioskowania, \(P\) oznacza
listę sekwentów-przesłanek, natomiast \(A\)
oznacza dodatkowe dane potrzebne w tych regułach, których zastosowanie nie jest
jednoznacznie wyznaczone przez same przesłanki. Przykładowo w regule \((\vee I_1)\) z samej przesłanki można odczytać jedynie
formułę \(\varphi\), ponieważ przesłanka ma postać \(\Gamma \vdash \varphi\).
Aby skonstruować konkluzję
\[
    \Gamma \vdash \varphi \vee \psi,
\]
system musi dodatkowo znać formułę \(\psi\). Nie jest ona wyznaczona przez
przesłankę, dlatego jest przekazywana jako dodatkowy argument zastosowania reguły.
Analogiczna sytuacja występuje w regule \((\vee I_2)\), gdzie dodatkowym argumentem
jest lewy składnik alternatywy.

Istotne jest to, że niepoprawne zastosowanie reguły nie jest reprezentowane przez
wartość logiczną \texttt{False}. Zamiast tego reguła zgłasza błąd. Jest to wygodne
z punktu widzenia dalszej konstrukcji systemu, ponieważ zastosowanie reguły można
traktować jako operację, która albo zwraca poprawnie skonstruowany sekwent, albo
w ogóle nie jest zdefiniowana dla podanych argumentów. Dzięki temu w pozostałych
częściach programu nie trzeba po każdym zastosowaniu reguły osobno sprawdzać,
czy otrzymany wynik jest poprawny. Błędne kandydatury mogą być po prostu
odrzucane przez przechwycenie błędu.

\section{Reprezentacja dowodów}
\label{sec:Reprezentacja dowodów}
W poprzednim rozdziale przyjęliśmy, że każda linia dowodu koduje jeden sekwent.
Na pierwszy rzut oka mogłoby się więc wydawać, że po sparsowaniu dowodu wystarczy
przechowywać wyłącznie sekwenty odpowiadające kolejnym liniom. Takie rozwiązanie
pozwalałoby odtworzyć, jakie sekwenty zostały kolejno otrzymane, ale traciłoby
informację o tym, jak poszczególne kroki zostały uzasadnione.

Sama reprezentacja sekwentu jest więc zbyt uboga jako reprezentacja linii dowodu.
Sekwent jest logicznym rezultatem danej linii, natomiast linia zawiera również dane
potrzebne do odtworzenia tego rezultatu: numer linii, formułę zapisaną jako jej tezę,
użytą regułę wnioskowania oraz indeksy występujące w uzasadnieniu. Bez tych
informacji sparsowany dowód byłby jedynie ciągiem otrzymanych sekwentów, a nie
strukturą pozwalającą odtworzyć zależności między krokami.

Sparsowana linia dowodu nie przechowuje oryginalnego tekstu linii, lecz informacje
wydobyte z tego tekstu oraz sekwent odtworzony na ich podstawie. Do informacji
odczytanych z zapisu należą: numer linii, lista indeksów występujących w jej
uzasadnieniu, obiekt reprezentujący zastosowaną regułę oraz formuła zapisana jako
teza linii. Na podstawie tych danych, a także wcześniej sparsowanej części dowodu,
konstruowany jest sekwent kodowany przez daną linię.


Można to przedstawić schematycznie w następujący sposób:
\[
\text{zapis linii}
\quad\longrightarrow\quad
(i,\varphi,R,J)
\quad\longrightarrow\quad
S_i.
\]
Tutaj \(i\) oznacza numer aktualnej linii, \(\varphi\) formułę zapisaną w tej linii,
\(R\) zastosowaną regułę, natomiast \(J=(j_1,\ldots,j_n)\) oznacza ciąg indeksów
podanych w uzasadnieniu linii. Indeksy te nie muszą być dokładnie listą przesłanek
reguły wnioskowania, ale w zależności od reguły mogą również zawierać indeksy linii z których czerpiemy dodatkowe argumenty do wywołania tej reguły, które nie są przesłankami. 
Na przykład linia postaci
\[
    \texttt{i. $\phi$ \quad   weakening, j, k}
\]
zawiera dwa indeksy, ale tylko jeden z nich wskazuje właściwą przesłankę reguły.
Sekwent \(S_k\) jest sekwentem, który ma zostać osłabiony, natomiast linia \(j\)
służy do odczytania formuły dodawanej do kontekstu. Jeśli więc
\(S_k = \Gamma \vdash \phi\), a tezą linii \(j\) jest \(\psi\), to system konstruuje
sekwent
\[
    \Gamma,\psi \vdash \phi.
\]
W tym przypadku dodatkowa formuła \(\psi\) nie pochodzi z formuły zapisanej
w aktualnej linii, lecz z tezy wcześniejszej linii wskazanej przez jeden z indeksów
w \(J\).

Innynym przykładem jest linia postaci
\[
    \texttt{i. $\phi \lor \psi \quad$ $\lor$-introduction1, j}
\]

Tutaj indeks \(j\) wskazuje właściwą przesłankę, czyli wcześniejszy sekwent
\(S_j = \Gamma \vdash \varphi\). Sama przesłanka nie wyznacza jednak formuły
\(\psi\), która ma zostać dodana jako prawa strona alternatywy. System odczytuje
ją więc z formuły zapisanej w aktualnej linii, czyli z prawego składnika
\(\varphi \vee \psi\), i przekazuje ją jako dodatkowy argument reguły.


Jeszcze innym przypadkiem są reguły, które nie odwołują się do żadnej wcześniejszej
linii. Przykładem jest linia postaci
\[
\texttt{1. $\phi$    assumption}
\]
W najprostszym przypadku odpowiada ona sekwentowi \(\phi \vdash \phi\). W naszym
systemie sekwent konstruowany dla takiej linii może jednak zawierać również niepusty
kontekst bazowy, obowiązujący w całym aktualnie rozważanym fragmencie dowodu.
Tego kontekstu nie da się odczytać z wcześniejszych sparsowanych linii, ponieważ
takowe jeszcze nie istnieją. Dlatego kontekst bazowy musi zostać przekazany explicite
jako dodatkowy argument podczas konstruowania linii.


Sparsowany dowód jest przechowywany jako słownik, który mapuje indeksy linii na
odpowiadające im sparsowane linie. Jeżeli przez \(D\) oznaczymy taki słownik, to
\(D[i]\) jest obiektem reprezentującym linię o numerze \(i\), a \(D[i].\mathrm{LP}\)
oznacza sekwent odtworzony dla tej linii. Dzięki temu podczas parsowania kolejnej
linii wystarczy odwołać się do indeksów zapisanych w jej uzasadnieniu i pobrać z
\(D\) potrzebne wcześniej skonstruowane obiekty.

Schematycznie parsowanie dowodu można więc opisać następująco:
\[
D_0 = \varnothing,
\]
\[
D_i = D_{i-1} \cup \{\, i \mapsto L_i \,\},
\]
gdzie \(L_i\) jest sparsowaną postacią \(i\)-tej linii dowodu, skonstruowaną przy
użyciu słownika \(D_{i-1}\). Oznacza to, że każda kolejna linia może korzystać
wyłącznie z informacji odtworzonych dla wcześniejszych linii. Taki sposób parsowania
odpowiada liniowej strukturze dowodu: późniejsze kroki mogą odwoływać się do
wcześniejszych, ale nie odwrotnie.

Po skonstruowaniu sparsowanej linii system sprawdza jeszcze, czy teza odtworzonego
sekwentu jest identyczna z formułą zapisaną w tekście linii. Ten krok jest konieczny,
ponieważ formuła występująca w zapisie linii jest traktowana jako deklaracja, a nie
jako samodzielne źródło informacji o całym sekwencie. Jeżeli reguła zastosowana do
odtworzonych przesłanek i argumentów prowadzi do innej tezy niż ta zapisana w
linii, linia zostaje odrzucona. W ten sposób parser pełni jednocześnie rolę
weryfikatora poprawności lokalnych kroków dowodu.

Niepoprawna linia nie jest oznaczana specjalną wartością logiczną, lecz powoduje
zgłoszenie błędu. Dotyczy to zarówno błędów składniowych, jak i sytuacji, w których
indeksy zapisane w uzasadnieniu nie pozwalają skonstruować poprawnego zastosowania
wskazanej reguły. Takie rozwiązanie upraszcza późniejsze etapy działania systemu.
Podczas generowania korpusu lub sprawdzania linii proponowanych przez model
językowy można po prostu próbować skonstruować sparsowaną linię; jeśli konstrukcja
się powiedzie, linia jest poprawna, a jeśli pojawi się błąd, kandydatura zostaje
odrzucona.

Ta sama reprezentacja jest wykorzystywana przy generowaniu korpusu. Program nie
tworzy linii treningowych przez ręczne sklejanie napisów, lecz najpierw konstruuje
poprawny obiekt linii, a dopiero potem wypisuje go w ustalonym formacie tekstowym.
Dzięki temu każda wygenerowana linia ma od razu odpowiadającą jej strukturę
logiczną, którą parser potrafi ponownie odtworzyć.


\subsection{Modele językowe}

W niniejszej pracy jednym z podstawowych narzędzi wykorzystywanych do konstrukcji
systemu dowodzącego są modele językowe. Modelem językowym nazywamy model, którego
zadaniem jest opisywanie regularności występujących w tekstach. W praktyce można
go rozumieć jako system, który dla danego początkowego fragmentu tekstu przewiduje,
jakie elementy mogą pojawić się dalej. W typowych zastosowaniach tekstem takim
może być tekst w języku naturalnym, fragment programu komputerowego albo odpowiedź
na pytanie. W tej pracy tekstem generowanym przez model będą natomiast dowody w
ustalonej notacji oraz zapisy rozbić problemów logicznych na prostsze podproblemy.

Tekst przekazywany do modelu nie jest zwykle analizowany jako ciąg znaków w
najprostszym sensie. Najpierw zostaje on podzielony na tokeny, czyli jednostki,
na których bezpośrednio operuje model. Tokenami mogą być pojedyncze znaki,
fragmenty słów, całe słowa albo inne fragmenty tekstu zależnie od przyjętej
metody kodowania. Po takim zakodowaniu tekst można traktować jako ciąg tokenów.
Zadaniem modelu językowego jest wtedy przewidywanie kolejnych tokenów na podstawie
tokenów wcześniejszych.

W pracy korzystam z modeli autoregresyjnych. Model autoregresyjny generuje tekst
krok po kroku. Mając pewien ciąg tokenów wejściowych, wyznacza rozkład
prawdopodobieństwa następnego tokenu. Następnie wybrany token jest dopisywany do
dotychczasowego tekstu i cały proces może zostać powtórzony. W ten sposób model,
który w pojedynczym kroku przewiduje tylko jeden następny token, może wygenerować
dłuższy fragment tekstu. Modele typu GPT są właśnie przykładami takich modeli
autoregresyjnych opartych na architekturze transformerów
\cite{vaswani2017attention,radford2019language}.

Trenowanie modelu językowego polega na uczeniu go przewidywania kolejnych tokenów
w tekstach pochodzących z korpusu treningowego. Model otrzymuje fragment tekstu
i ma przewidzieć tokeny występujące po nim. Następnie jego parametry są zmieniane
tak, aby zwiększać prawdopodobieństwo tych tokenów, które rzeczywiście wystąpiły
w danych treningowych. Można więc powiedzieć, że model uczy się regularności
obecnych w korpusie: zarówno prostych zależności lokalnych między sąsiednimi
tokenami, jak i bardziej złożonych wzorców pojawiających się w dłuższych
fragmentach tekstu.

W niniejszej pracy szczególnie istotne jest to, że tekstem, którego regularności
ma nauczyć się model, nie jest zwykły tekst w języku naturalnym, lecz zapis
dowodu albo zapis rozbicia problemu logicznego. Dowody w rozważanej notacji mają
bardzo regularną strukturę: każda linia zawiera formułę, nazwę zastosowanej
reguły oraz indeksy wcześniejszych linii, do których dana reguła się odwołuje.
Podobnie zapisy rozbić problemów na podproblemy mają ustalony format tekstowy.
Pozwala to potraktować generowanie kolejnych kroków dowodu jako szczególny
przypadek generowania tekstu.

Jako bazę implementacyjną wykorzystuję projekt \texttt{nanoGPT}
\cite{karpathy2022nanogpt}. Jest to niewielka implementacja modelu typu GPT,
która pozwala trenować modele językowe na własnych korpusach tekstowych. W moim
systemie trenowane są modele wyspecjalizowane do konkretnych zadań związanych
z dowodzeniem. Jeden z nich służy do generowania dowodów sekwentów z założeniami,
natomiast drugi do generowania rozbić problemów na prostsze podproblemy. Oba
modele uczone są na specjalnie przygotowanych korpusach, których konstrukcja
zostanie omówiona w dalszej części pracy.

Należy jednak podkreślić, że model językowy nie jest w tej pracy traktowany jako
samodzielny system dowodzący. Sam fakt, że model wygenerował tekst przypominający
dowód, nie oznacza jeszcze, że jest to dowód poprawny formalnie. Model może
wskazać niewłaściwą regułę, odwołać się do niepasujących linii albo wygenerować
formułę, która nie jest konkluzją wynikającą z podanych przesłanek. Dlatego model
pełni rolę komponentu heurystycznego: proponuje kolejne kroki dowodu albo możliwe
rozbicia problemu, natomiast ich poprawność jest sprawdzana przez część symboliczną
systemu.

Takie połączenie modelu językowego i formalnego weryfikatora jest naturalne dla
rozważanego problemu. Z jednej strony klasyczne, czysto symboliczne przeszukiwanie
przestrzeni dowodów może być trudne, ponieważ często nie wiadomo, którą regułę
wnioskowania zastosować i jakie formuły pomocnicze należy wprowadzić. Z drugiej
strony model językowy może nauczyć się z korpusu, jakie kroki dowodowe często
pojawiają się w podobnych sytuacjach. Nie gwarantuje on jednak poprawności tych
kroków. Weryfikator pełni więc rolę filtra, który odrzuca propozycje niepoprawne
i pozwala wykorzystać tylko te fragmenty wygenerowanego tekstu, które rzeczywiście
kodują poprawne zastosowania reguł wnioskowania.

W rezultacie system nie polega wyłącznie na generowaniu tekstu przez model, ale
łączy dwa podejścia. Część neuronowa służy do podpowiadania potencjalnych kierunków
dowodu, natomiast część symboliczna odpowiada za jednoznaczną interpretację
wygenerowanych linii i sprawdzenie ich poprawności. Dzięki temu można wykorzystać
zdolność modeli językowych do rozpoznawania regularności w danych tekstowych,
jednocześnie zachowując formalną kontrolę nad poprawnością konstruowanych dowodów.






\chapter{Strategie konstruowania dowodów}

\section{Motywacja}
Celem systemu nie jest jedynie wygenerowanie tekstu przypominającego dowód
w dedukcji naturalnej, lecz znalezienie poprawnego dowodu danej tautologii.
Model językowy pełni w tym procesie rolę generatora kandydatów: proponuje kolejne
linie dowodu albo inne kroki pomocnicze, ale nie decyduje samodzielnie o ich
poprawności. Każdy wygenerowany fragment musi zostać zweryfikowany przez opisany
wcześniej parser i mechanizm sprawdzania reguł wnioskowania.

Takie rozdzielenie ról jest kluczowe. Model językowy może dobrze naśladować składnię
dowodów i często proponować sensowne kroki, ale sam z siebie nie gwarantuje
poprawności logicznej. Może wygenerować linię o poprawnym wyglądzie, która jednak
odwołuje się do niewłaściwych wcześniejszych linii, używa reguły w niepoprawny
sposób albo wprowadza formułę, której nie da się uzasadnić w danym kontekście.
Dlatego w systemie model nie jest traktowany jako samodzielny system, lecz jako źródło
propozycji filtrowanych przez formalny weryfikator.

Najprostsze podejście polega na liniowym dopisywaniu kolejnych linii dowodu.
System przekazuje modelowi cel dowodowy oraz dotychczas wygenerowany fragment
dowodu, a następnie próbuje dopisać jedną nową linię. Jeżeli po jej dopisaniu dowód
nadal poprawnie się parsuje, linia zostaje zaakceptowana; w przeciwnym razie jest
odrzucana. Proces generacji trwa do momentu uzyskania linii kodującej poszukiwany
sekwent albo do chwili, w której dalsze rozszerzanie dowodu zostaje uznane za
niecelowe. Tak powstaje pierwsza strategia, w której dowód konstruowany jest jako
jedna rosnąca sekwencja lokalnie poprawnych kroków.

Podejście to ma jednak istotne ograniczenie. Konstruowanie dowodu rzadko jest
procesem całkowicie liniowym. Nawet jeśli każda kolejna linia jest poprawna, może
okazać się, że wybrana ścieżka nie prowadzi do celu albo prowadzi do niego bardzo
okrężną drogą. W praktyce potrzebne jest więc pewne przeszukiwanie przestrzeni
możliwych częściowych dowodów: system powinien móc rozwijać kilka alternatywnych
prób i wracać do bardziej obiecujących fragmentów. Tę ideę realizuje druga
strategia, która przechowuje wiele częściowych dowodów i wybiera spośród nich te,
które warto dalej rozwijać.

Dla trudniejszych formuł nawet takie lokalne przeszukiwanie może być niewystarczające.
Problemem nie jest wtedy tylko wybór następnej linii, lecz wybór ogólnego planu
dowodu. Często łatwiej jest najpierw zdecydować, jaką regułą można byłoby otrzymać
cel, a następnie spróbować udowodnić jej przesłanki jako osobne podproblemy. W ten
sposób otrzymujemy trzecią strategię, w której system może rozbić początkowy
problem na prostsze sekwenty, a następnie szukać dowodów dla nich. Strategia ta
łączy generowanie dowodów w notacji liniowej z przeszukiwaniem drzewa problemów
i podproblemów.

Kolejne strategie powstały więc jako odpowiedź na ograniczenia poprzednich.
Pierwsza strategia sprawdza, czy model potrafi lokalnie dopisywać poprawne linie
dowodu. Druga dodaje mechanizm wyboru między wieloma częściowymi próbami,
ale nadal wykorzystuje pierwszą strategię jako procedurę rozwijania pojedynczego
fragmentu dowodu. Trzecia zmienia poziom poszukiwania: oprócz dopisywania linii
pozwala planować dowód przez dekompozycję celu na podproblemy. Również ona nie
zastępuje wcześniejszych metod, lecz korzysta z nich jako komponentów pomocniczych:
po rozbiciu celu na podproblemy próbujemy znaleźć dowody tych podproblemów przy 
użyciu drugiej strategii.

W ten sposób strategie tworzą hierarchię. Każda kolejna strategia rozszerza
poprzednią, traktując ją jako niezależny komponent realizujący prostsze zadanie. Dzięki
temu system można opisywać warstwowo: najpierw jako generator pojedynczej
poprawnej kontynuacji dowodu, następnie jako mechanizm przeszukiwania wielu
takich kontynuacji, a ostatecznie jako procedurę rozbijania problemu dowodowego
na mniejsze problemy, dla których ponownie uruchamiane są wcześniejsze mechanizmy.
W dalszej części rozdziału strategie te zostaną opisane w tej kolejności.



\section{Strategia pierwsza: liniowe dopisywanie kroków dowodu}

\subsection{Idea strategii}

Pierwsza strategia jest najprostszym sposobem wykorzystania modelu językowego
w systemie dowodzącym. Jej celem jest sprawdzenie, czy model potrafi generować
kolejne poprawne kroki dowodu w notacji liniowej, jeżeli otrzyma opis celu oraz
dotychczasowy fragment dowodu. Strategia ta nie próbuje planować całego dowodu
z góry ani nie rozważa wielu alternatywnych ścieżek. Dowód jest konstruowany
sekwencyjnie: w każdym kroku system próbuje dopisać jedną nową linię do końca
aktualnego fragmentu.

W tym sensie strategia pierwsza traktuje model jako lokalny generator kontynuacji.
Wejściem dla modelu jest prompt kodujący sekwent, który ma zostać udowodniony,
oraz ewentualnie dotychczas wygenerowane linie dowodu. Wyjściem nie jest cały
dowód, lecz kandydatura na następną linię. Poprawność tej kandydatury nie jest
oceniana przez model, lecz przez formalny parser i mechanizm sprawdzania reguł
wnioskowania.

Schematycznie można to przedstawić następująco:
\[
\begin{array}{c}
\text{częściowy}\\
\text{dowód}
\end{array}
\longrightarrow
\begin{array}{c}
\text{propozycja}\\
\text{następnej linii}
\end{array}
\longrightarrow
\text{weryfikacja}
\longrightarrow
\begin{array}{c}
\text{zaakceptowanie}\\
\text{lub odrzucenie}
\end{array}
\]

W implementacji tej strategii odpowiada podstawowa procedura lokalnego rozwijania
częściowego dowodu, oznaczona jako \texttt{silliest\_generate\_proof}. Nazwa ta
odzwierciedla prototypowy charakter pierwszej wersji algorytmu: strategia jest
celowo prosta i służy jako punkt wyjścia dla bardziej złożonych metod opisanych
w dalszej części rozdziału.


\subsection{Generowanie i weryfikacja pojedynczej linii}

Podstawowa operacja wykonywana przez strategię polega na próbie dopisania jednej
nowej linii do aktualnego fragmentu dowodu. Nie odbywa się to jednak przez zupełnie
swobodne wygenerowanie dowolnego ciągu znaków przez model. System narzuca
częściową strukturę generowanej linii, tak aby model odpowiadał przede wszystkim
za wybór formuły i uzasadnienia, a nie za odtwarzanie elementów czysto technicznych,
takich jak numeracja linii.

Najpierw analizowany jest aktualny prompt. Jeżeli zawiera on wyłącznie cel dowodowy,
następna linia otrzymuje numer \(1\). Jeżeli w prompcie znajdują się już wcześniejsze
linie dowodu, system odczytuje numer ostatniej z nich i rozpoczyna nową linię od
kolejnego numeru. Dzięki temu model nie musi samodzielnie utrzymywać poprawnej
numeracji dowodu, co zmniejsza liczbę błędów czysto składniowych.

Następnie generowana jest formuła zapisana w nowej linii. Na tym etapie system
ogranicza zbiór akceptowanych tokenów. Dopuszczalne są symbole logiczne, nawiasy,
stałe \(\top\) i \(\bot\), spacje oraz zmienne występujące w dowodzonym sekwencie.
Jeżeli dowód ma niepusty kontekst bazowy, zbiór dopuszczalnych zmiennych jest
wyznaczany na podstawie zarówno tezy, jak i formuł występujących w kontekście.
Celem tego ograniczenia jest uniknięcie sytuacji, w której model wprowadza do
dowodu nowe, przypadkowe zmienne, niezwiązane z rozważanym problemem.

Generowanie formuły trwa do momentu, w którym pojawi się separator oddzielający
formułę od nazwy reguły uzasadniającej linię. W przyjętym formacie rolę tę pełnią cztery kolejne
spacje. Po ich wygenerowaniu system przechodzi do drugiej części linii, czyli do
uzasadnienia zawierającego nazwę reguły oraz ewentualne indeksy wcześniejszych
linii. Ta część jest generowana już jako kontynuacja rozpoczętej linii i kończy się
w momencie wygenerowania znaku nowej linii.

Można to opisać schematycznie następująco:
\[
\text{prompt}
\longrightarrow
\text{numer nowej linii}
\longrightarrow
\text{formuła}
\longrightarrow
\text{uzasadnienie}.
\]
W implementacji ten etap odpowiada procedurze generowania następnej linii, która
na podstawie aktualnego prompta tworzy kandydaturę na kolejną linię dowodu.

Po wygenerowaniu kandydackiej linii system dopisuje ją tymczasowo do aktualnego
fragmentu dowodu i próbuje ponownie sparsować całość. Sprawdzenie to nie ogranicza
się do kontroli składni. Parser próbuje odtworzyć strukturę każdej linii,
skonstruować odpowiadające jej sekwenty i zweryfikować, czy zastosowane reguły
rzeczywiście prowadzą do zapisanych tez. Jeżeli parsowanie kończy się błędem,
wygenerowana linia zostaje odrzucona.

Kandydacka linia może zostać odrzucona z kilku powodów. Po pierwsze, może nie
spełniać wymagań składniowych przyjętego formatu. Po drugie, może zawierać
formułę niezgodną z ograniczeniami na zmienne występujące w dowodzonym problemie
(sprawdzenie tego warunku jest defacto nadmiarowe, ponieważ jest on sprawdzany już na etapie generowania linii).
Po trzecie, może wyglądać jak poprawna linia dowodu, ale odwoływać się do
wcześniejszych linii w sposób niezgodny ze schematem wskazanej reguły. W każdym
z tych przypadków parser lub weryfikator zgłasza błąd, a wygenerowana kandydatura
nie jest dopisywana do dowodu.

Jeżeli natomiast dopisanie linii nie powoduje błędu parsowania, linia zostaje
zaakceptowana. Oznacza to, że udało się odtworzyć sekwent kodowany przez tę linię,
a jego teza zgadza się z formułą zapisaną w tekście. Po zaakceptowaniu linii system
traktuje ją tak samo jak wcześniejsze linie: może ona zostać wykorzystana jako
uzasadnienie kolejnych kroków dowodu.

W praktyce strategia działa więc jako pętla generowania i filtrowania. Jeśli przez
\(D_k\) oznaczymy aktualny częściowy dowód, a przez \(\ell\) kandydaturę na nową
linię, to po jej tymczasowym dopisaniu otrzymujemy
\[
D_{k+1}^{\ast}=D_k+\ell.
\]
Następnie system wykonuje weryfikację:
\[
D_{k+1}=
\begin{cases}
D_{k+1}^{\ast}, & \text{jeśli } D_{k+1}^{\ast} \text{ jest poprawny},\\
D_k, & \text{w przeciwnym przypadku.}
\end{cases}
\]
Takie rozwiązanie pozwala oddzielić zadanie generowania od zadania sprawdzania.
Model językowy odpowiada za proponowanie potencjalnie sensownych kontynuacji,
natomiast poprawność logiczna pozostaje kontrolowana przez formalny mechanizm
opisany w poprzednim rozdziale.
\subsection{Warunki zakończenia}

Strategia kończy działanie sukcesem wtedy, gdy ostatnia zaakceptowana linia koduje
sekwent docelowy. Innymi słowy, po sparsowaniu aktualnego fragmentu dowodu system
sprawdza, czy sekwent odtworzony dla ostatniej linii jest równy sekwentowi, który
miał zostać udowodniony. Jeśli tak jest, wygenerowany tekst stanowi pełny dowód
celu.

Działanie strategii może zakończyć się także bez znalezienia dowodu. Pierwszym
ograniczeniem jest maksymalna liczba nowych linii, które wolno dopisać do prompta.
Chroni to system przed generowaniem bardzo długich fragmentów, które lokalnie
pozostają poprawne, ale nie zbliżają się do celu. Drugim ograniczeniem jest maksymalna
liczba kolejnych nieudanych prób dopisania linii. Jeżeli model przez wiele prób
z rzędu generuje kandydatury odrzucane przez parser, system uznaje, że dalsze
kontynuowanie tej samej ścieżki jest mało obiecujące.

W przypadku zakończenia bez sukcesu strategia może mimo to zwrócić częściowy
dowód, czyli prompt z dopisanymi poprawnymi liniami. Taki fragment nie musi jeszcze
dowodzić celu, ale może stanowić postęp względem punktu wyjścia. Właśnie dlatego
pierwsza strategia może być później użyta jako komponent w strategii drugiej:
nie musi samodzielnie zawsze znaleźć pełnego dowodu, wystarczy, że potrafi lokalnie
rozwijać częściowy dowód o kolejne poprawne kroki.

\subsection{Ograniczenia strategii}

Największym ograniczeniem pierwszej strategii jest jej liniowy charakter. System
rozwija tylko jedną ścieżkę dowodu i nie przechowuje alternatywnych kontynuacji.
Jeżeli model wybierze poprawną, ale mało użyteczną linię, strategia nie ma mechanizmu
powrotu do wcześniejszego stanu i wyboru innej możliwości. Może więc utknąć w
ciągu kroków, które są lokalnie poprawne, ale nie prowadzą do dowodu celu.

Drugim ograniczeniem jest brak globalnego planowania. Strategia nie analizuje
struktury celu w taki sposób, aby zdecydować, jakie podcele warto najpierw osiągnąć.
Każdy krok jest generowany na podstawie dotychczasowego tekstu dowodu, ale system
nie rozważa jawnie, jaką regułą można byłoby zakończyć dowód ani jakie przesłanki
należałoby wcześniej udowodnić. W prostych przypadkach lokalne dopisywanie linii
może wystarczyć, ale dla bardziej złożonych formuł często potrzebne jest planowanie
na wyższym poziomie.

Skuteczność tej strategii zależy również od parametrów generacji, takich jak
temperatura oraz liczba dopuszczalnych prób wygenerowania kolejnej linii. Nie jest
to jednak osobny problem konstrukcyjny, lecz ogólna konsekwencja użycia modelu
językowego jako generatora kandydatów.

Mimo tych ograniczeń strategia pierwsza jest ważnym elementem całego systemu.
Stanowi podstawowy mechanizm lokalnego rozszerzania dowodu, który może być
używany jako komponent w bardziej złożonych metodach. Druga strategia zachowuje
ten mechanizm generowania pojedynczych kroków, ale dodaje możliwość rozwijania
wielu częściowych dowodów równolegle, co częściowo rozwiązuje problem utknięcia
w jednej niekorzystnej ścieżce.


\section{Strategia druga: przeszukiwanie wielu częściowych dowodów}

\subsection{Motywacja}

Pierwsza strategia konstruuje tylko jeden częściowy dowód. Oznacza to, że każda
zaakceptowana linia staje się trwałą częścią aktualnej próby. Jeżeli później okaże się,
że dana kontynuacja nie prowadzi w stronę rozwiązania, system nie ma możliwości
powrotu do wcześniejszego etapu i wyboru innej ścieżki rozumowania. W praktyce
jest to poważne ograniczenie, ponieważ poprawność lokalna pojedynczej linii nie
gwarantuje jeszcze, że linia ta będzie użyteczna dla całego dowodu.

Druga strategia próbuje ograniczyć ten problem przez jednoczesne przechowywanie
wielu częściowych dowodów. Zamiast rozwijać wyłącznie jedną sekwencję linii, system
utrzymuje zbiór kandydatów. W każdej iteracji wybierany jest jeden częściowy dowód,
a następnie podejmowana jest próba dopisania do niego kolejnej poprawnej linii.
Jeżeli powstały w ten sposób fragment jest poprawny i nie wystąpił wcześniej, zostaje
dodany do zbioru kandydatów.

Można więc powiedzieć, że druga strategia zmienia charakter generowania dowodu
z liniowego dopisywania kolejnych kroków na przeszukiwanie przestrzeni częściowych
dowodów. Pojedynczy nieoptymalny wybór nie przekreśla całego procesu, ponieważ
inne częściowe dowody pozostają dostępne i mogą zostać rozwinięte w późniejszych
iteracjach.

Schematycznie można to przedstawić następująco:
\[
\small
\begin{array}{c}
\text{zbiór częściowych}\\[-0.2ex]
\text{dowodów}
\end{array}
\;\longrightarrow\;
\begin{array}{c}
\text{wybór jednego}\\[-0.2ex]
\text{kandydata}
\end{array}
\;\longrightarrow\;
\begin{array}{c}
\text{próba dopisania}\\[-0.2ex]
\text{nowej linii}
\end{array}
\;\longrightarrow\;
\begin{array}{c}
\text{aktualizacja zbioru}\\[-0.2ex]
\text{kandydatów}
\end{array}.
\]

\subsection{Zbiór częściowych dowodów}

Podstawową strukturą wykorzystywaną przez drugą strategię jest zbiór częściowych
dowodów. Na początku zawiera on jedynie prompt opisujący sekwent, który należy
udowodnić. W kolejnych iteracjach do zbioru dodawane są nowe częściowe dowody
powstałe przez dopisanie poprawnej linii do jednego z wcześniej istniejących
kandydatów.

Każdy element tego zbioru jest tekstem zawierającym opis celu dowodowego oraz
dotychczas wygenerowany fragment dowodu. Dzięki temu może zostać bezpośrednio
podany modelowi językowemu jako kontekst generacji. 

W każdej iteracji system wykonuje następujące kroki:
\begin{algorithm}[H]
\caption{Rozwijanie zbioru częściowych dowodów}
\begin{algorithmic}[1]
\State Wybierz częściowy dowód \(D\) ze zbioru kandydatów \(P\).
\State Spróbuj dopisać do \(D\) jedną nową poprawną linię, otrzymując \(D'\).
\State Sprawdź poprawność otrzymanego fragmentu \(D'\).
\If{\(D'\) kończy dowód celu}
    \State \Return \(D'\)
\ElsIf{\(D'\) jest poprawnym nowym częściowym dowodem}
    \State Dodaj \(D'\) do zbioru kandydatów \(P\).
\EndIf
\end{algorithmic}
\end{algorithm}

Ważne jest, że druga strategia nie zastępuje mechanizmu generowania pojedynczej
linii, lecz używa go jako podprocedury. Próba rozwinięcia częściowego dowodu nadal
polega na tym, że model proponuje kandydacką linię, a formalny parser i weryfikator
sprawdzają jej poprawność. Różnica polega na tym, że poprawne rozwinięcia nie
muszą tworzyć jednej liniowej historii. Mogą współistnieć jako alternatywne ścieżki
poszukiwania dowodu.

\subsection{Ocena prawdopodobieństwa częściowego dowodu}

Jeżeli zbiór częściowych dowodów zawiera wiele kandydatów, trzeba zdecydować,
który z nich rozwijać w następnej iteracji. Najprostsze rozwiązanie polegałoby na
losowaniu jednostajnym, jednak ignorowałoby ono to, że nie wszystkie częściowe dowody są tak samo prawdopodobne.
Model nie tylko generuje tekst, ale również przypisuje kolejnym tokenom
prawdopodobieństwa. Można więc użyć tych prawdopodobieństw jako przybliżonej
oceny tego, jak naturalny dla modelu jest dany częściowy dowód.

Dla każdego częściowego dowodu obliczana jest wartość zależna od prawdopodobieństw
tokenów tworzących wygenerowaną część dowodu. Ponieważ iloczyn wielu
prawdopodobieństw bardzo szybko staje się liczbą bliską zeru, wygodniej sumować
logarytmy prawdopodobieństw. Aby nie faworyzować nadmiernie krótkich fragmentów
tylko dlatego, że składają się z mniejszej liczby tokenów, suma ta jest normalizowana
przez długość analizowanego tekstu.

Jeżeli częściowy dowód oznaczymy przez \(D\), a jego kolejne tokeny przez
\(t_1,\ldots,t_m\), to jego ocenę można zapisać schematycznie jako
\[
    score(D)
    =
    \frac{1}{m}
    \sum_{i=1}^{m}
    \log P(t_i \mid t_1,\ldots,t_{i-1}, \text{prompt celu}).
\]

Wartość \(score(D)\) jest jednak nieokreślona, gdy \(m=0\), czyli dla początkowego
prompta, który nie zawiera jeszcze żadnych wygenerowanych tokenów. W takim przypadku
najpierw wyliczamy \(score\) dla pozostałych częściowych dowodów, a następnie
ustawiamy \(score\) dla pustego prompta na wartość średnią z tych ocen. Intuicja
jest taka, że początkowy prompt nie powinien być ani stale faworyzowany, ani całkowicie
pomijany. Z jednej strony nie zawiera on jeszcze żadnego postępu w konstrukcji
dowodu, więc zbyt częste wybieranie go prowadziłoby do wielokrotnego rozpoczynania
generacji od początku. Z drugiej strony jest to kandydat najbardziej neutralny:
nie zawiera żadnych wcześniejszych decyzji, które mogłyby skierować dowód w mało
użyteczną stronę. Przypisanie mu średniej oceny pozwala więc zachować możliwość
powrotu do początku rozumowania, ale nie czyni z niego domyślnie najatrakcyjniejszej
ścieżki.


Kolejnym problemem jest to, że wartości \(score(D)\), jako średnie logarytmów
prawdopodobieństw, są w praktyce niedodatnie. Mogą być równe \(0\) jedynie w
granicznym przypadku, gdy każdy kolejny token ma prawdopodobieństwo warunkowe
równe \(1\). Ponadto wartości te nie sumują się do \(1\), a więc nie mogą być
bezpośrednio interpretowane jako rozkład prawdopodobieństwa. Aby przekształcić je
w rozkład na zbiorze kandydatów, używamy funkcji softmax:
\[
    P(D_i)
    =
    \frac{\exp(score(D_i))}
    {\sum_j \exp(score(D_j))}.
\]
W ten sposób każdemu częściowemu dowodowi przypisana zostaje dodatnia wartość,
a suma wszystkich takich wartości wynosi \(1\). Kandydaci o wyższym wyniku
\(score\), czyli bardziej prawdopodobni według modelu, otrzymują większe
prawdopodobieństwo wyboru w następnej iteracji.


\subsection{Pełny przebieg strategii}

Po wprowadzeniu zbioru częściowych dowodów oraz sposobu przypisywania im
prawdopodobieństw wyboru można przedstawić całą drugą strategię jako jedną pętlę
przeszukiwania. W każdej iteracji system aktualizuje oceny kandydatów, wybiera jeden
z częściowych dowodów zgodnie z rozkładem otrzymanym po zastosowaniu funkcji
softmax, a następnie próbuje rozwinąć go o jedną poprawną linię.

\begin{algorithm}[H]
\caption{Przeszukiwanie wielu częściowych dowodów}
\begin{algorithmic}[1]
\State \(P \gets \{\text{prompt opisujący cel } G\}\)
\While{nie przekroczono limitu iteracji}
    \State Oblicz wartości \(score(D)\) dla kandydatów \(D \in P\).
    \State Przekształć wartości \(score(D)\) w rozkład prawdopodobieństwa na \(P\).
    \State Wylosuj częściowy dowód \(D \in P\) zgodnie z otrzymanym rozkładem.
    \State Spróbuj dopisać do \(D\) jedną poprawną linię, otrzymując \(D'\).
    \If{\(D'\) jest pełnym dowodem celu \(G\)}
        \State \Return \(D'\)
    \ElsIf{\(D'\) jest poprawnym nowym częściowym dowodem}
        \State \(P \gets P \cup \{D'\}\)
    \EndIf
\EndWhile
\State \Return informacja o niepowodzeniu
\end{algorithmic}
\end{algorithm}

W powyższym schemacie etap dopisania nowej linii korzysta z procedury opisanej
przy pierwszej strategii. Druga strategia nie zmienia więc sposobu weryfikacji
pojedynczego kroku dowodu, lecz zmienia organizację całego procesu poszukiwania:
zamiast jednej rozwijanej sekwencji utrzymywany jest zbiór alternatywnych prób,
spośród których kolejne są wybierane probabilistycznie.

Niepowodzenie algorytmu oznacza jedynie, że dowód nie został znaleziony w ramach
przyjętych ograniczeń obliczeniowych. Nie jest to informacja o niedowodliwości
danego sekwentu.

\subsection{Ograniczenia strategii}

Druga strategia usuwa najważniejsze ograniczenie strategii liniowej: system nie jest
już przywiązany do jednej ścieżki generacji. Może utrzymywać wiele alternatywnych
częściowych dowodów i stopniowo rozwijać te, które wydają się bardziej obiecujące.
Nie rozwiązuje to jednak wszystkich problemów związanych z automatycznym
konstruowaniem dowodów.

Po pierwsze, strategia nadal ma charakter lokalny. Rozwijanie częściowego dowodu
polega na dopisaniu kolejnej linii, a nie na świadomym zaplanowaniu całej struktury
dowodu. System może więc generować poprawne, ale mało użyteczne kroki, które
zwiększają długość dowodu bez realnego przybliżania go do celu.

Po drugie, ocena prawdopodobieństwa częściowego dowodu nie jest tym samym co
ocena jego matematycznej przydatności. Model może wysoko oceniać fragmenty
podobne do danych treningowych, ale nie oznacza to jeszcze, że są one najlepszą
drogą do udowodnienia konkretnej tezy. Prawdopodobieństwo generacji jest więc
heurystyką, a nie kryterium poprawności ani gwarancją sukcesu.

Po trzecie, liczba możliwych częściowych dowodów rośnie bardzo szybko. Każda
zaakceptowana linia może prowadzić do kolejnych rozwinięć, a wiele z nich może być
formalnie poprawnych. W rezultacie strategia wymaga limitów liczby iteracji oraz
mechanizmów wyboru kandydatów do rozwijania. Limity te są konieczne obliczeniowo,
ale sprawiają, że przeszukiwanie pozostaje niepełne.

Najważniejsze ograniczenie polega jednak na tym, że druga strategia nadal nie
wykonuje dekompozycji celu na podcele. System próbuje znaleźć dowód przez
rozszerzanie częściowych tekstowych dowodów, ale nie ma osobnego etapu, w którym
analizuje strukturę dowodzonego sekwentu i wybiera regułę pozwalającą rozbić go
na prostsze problemy. W trudniejszych przypadkach właśnie taka dekompozycja może
być bardziej naturalna niż dopisywanie kolejnych linii od początku dowodu.

To prowadzi do trzeciej strategii, w której system może najpierw rozbić dowodzony
sekwent na listę prostszych podproblemów, a następnie próbować udowodnić każdy
z nich osobno.





\section{Strategia trzecia: dekompozycja celu i drzewo poszukiwania}


\subsection{Motywacja}

Druga strategia pozwala prowadzić wiele częściowych dowodów równolegle, dzięki
czemu system nie jest przywiązany do jednej liniowej ścieżki generacji. Nadal jednak
rozwijanie dowodu odbywa się przez dopisywanie kolejnych linii do już istniejącego
fragmentu. System wybiera więc pewien częściowy dowód i próbuje znaleźć jego
następny poprawny krok. Nie analizuje natomiast osobno, jaka ogólna struktura
dowodu byłaby potrzebna, aby otrzymać zadany sekwent.

Taki sposób działania odpowiada naturalnemu zapisowi dowodu w notacji Fitcha:
zaczynamy od założeń, aksjomatów lub wcześniej uzyskanych sekwentów, a następnie
stosujemy reguły wnioskowania, aby dojść do tezy. Jest to konstruowanie dowodu
w kierunku od przesłanek do celu. W wielu sytuacjach wygodny jest jednak kierunek
przeciwny. Zamiast pytać, jaką kolejną linię można dopisać do częściowego dowodu,
można zacząć od sekwentu, który chcemy udowodnić, i zapytać, jakie prostsze
sekwenty wystarczyłoby udowodnić, aby otrzymać go przez zastosowanie jednej
z reguł wnioskowania.

Można więc patrzeć na konstrukcję dowodu z dwóch stron. Pierwsza strategia buduje
dowód od wcześniejszych linii ku końcowej tezie. Trzecia strategia zaczyna od celu
i próbuje rozłożyć go na podcele, które następnie mogą być dowodzone osobno.
W notacji drzewiastej odpowiada to budowaniu drzewa dowodu od korzenia ku
liściom: w korzeniu znajduje się sekwent, który chcemy udowodnić, a kolejne poziomy
drzewa zawierają przesłanki reguł, których udowodnienie wystarcza do otrzymania
sekwentów znajdujących się wyżej.

Przykładowo, jeżeli celem jest udowodnienie koniunkcji, naturalnym krokiem jest
próba udowodnienia osobno jej dwóch składników. Jeżeli celem jest udowodnienie
implikacji, można przyjąć jej poprzednik jako dodatkowe założenie i następnie
spróbować udowodnić jej następnik. W obu przypadkach nie zaczynamy od pytania
o pierwszą linię dowodu, lecz od pytania o to, jakie podproblemy prowadzą do
rozwiązania problemu wejściowego.

Trzecia strategia rozszerza więc wcześniejsze podejścia o możliwość dekompozycji
celu. System może najpierw rozbić dowodzony sekwent na listę prostszych
podproblemów, a dopiero potem próbować udowodnić każdy z nich za pomocą
wcześniejszych metod generowania dowodów. Jeżeli wszystkie podproblemy danego
rozbicia zostaną rozwiązane, to rozwiązany zostaje również problem wejściowy.

\subsection{Rozbicie problemu na podproblemy}

Rozbiciem problemu na podproblemy nazywamy wskazanie takiej reguły
wnioskowania oraz takiej listy sekwentów-przesłanek, że sekwent wejściowy
można otrzymać przez zastosowanie tej reguły do tych przesłanek.

Aby zapisać to precyzyjniej, dla każdej reguły \(R\) oraz listy sekwentów
\(S_1,\ldots,S_n\) wprowadzamy oznaczenie
\[
R(S_1,\ldots,S_n),
\]
które oznacza zbiór wszystkich sekwentów mogących wystąpić jako konkluzje
poprawnego zastosowania reguły \(R\) do przesłanek \(S_1,\ldots,S_n\).


Jeżeli więc chcemy udowodnić sekwent \(S\), to szukamy reguły \(R\) oraz
sekwentów \(S_1,\ldots,S_n\), dla których zachodzi
\[
S \in R(S_1,\ldots,S_n).
\]
Wtedy zadanie udowodnienia \(S\) można zastąpić zadaniami udowodnienia
\(S_1,\ldots,S_n\).

Takie rozbicie można interpretować jako jeden krok budowania drzewa dowodu
w kierunku przeciwnym do kierunku zapisu dowodu w notacji Fitcha. W notacji
Fitcha zaczynamy zwykle od prostych przesłanek i dopisujemy kolejne linie aż
do uzyskania celu. W dekompozycji celu zaczynamy natomiast od sekwentu,
który chcemy otrzymać, i pytamy, jakie przesłanki mogłyby prowadzić do niego
przez jedną z reguł wnioskowania.

Schematycznie:
\[
\small
\begin{array}{c}
\text{problem}\\[-0.2ex]
S
\end{array}
\;\longrightarrow\;
\begin{array}{c}
\text{wybór}\\[-0.2ex]
\text{reguły } R
\end{array}
\;\longrightarrow\;
\begin{array}{c}
\text{podproblemy}\\[-0.2ex]
S_1,\ldots,S_n
\end{array}
\]
przy czym warunkiem poprawności rozbicia jest
\[
S \in R(S_1,\ldots,S_n).
\]

Rozbicie z założenia powinny generować podproblemy o dowodach prostszych niż dowód problemu początkowego.
Co więcej trzecia strategia nie zastępuje generowania dowodów liniowych, ale
organizuje je w większą strukturę: zamiast od razu znaleźć jeden
pełny dowód sekwentu, system może rozbić go na mniejsze sekwenty i dla
każdego z nich uruchomić szukanie dowodu w formie liniowej. Szukanie to odbywa się
przy pomocy strategii drugiej.


\subsection{Format zapytania o rozbicie}

W trzeciej strategii model językowy jest używany nie tylko do generowania kolejnych
linii dowodu, lecz także do proponowania rozbić celu na podproblemy. Jest to osobne
zadanie generatywne. W przypadku dowodzenia liniowego model otrzymuje zapis celu
oraz dotychczasowy fragment dowodu i ma zaproponować jego dalszy ciąg. W przypadku
rozbijania problemu model otrzymuje sekwent, który chcemy udowodnić, i ma wskazać
regułę wnioskowania oraz listę podproblemów, których udowodnienie pozwoli otrzymać
sekwent wejściowy.

Z tego powodu w systemie wykorzystywane są dwa modele językowe o tej samej
architekturze transformera, ale trenowane na różnych typach danych. Pierwszy model
uczy się generować dowody w notacji tekstowej, czyli dopisywać kolejne linie
częściowego dowodu. Drugi model uczy się generować rozbicia problemów na
podproblemy. Oba modele korzystają z podobnej idei: proponują tekstowy kandydat,
który następnie jest sprawdzany przez formalny parser i mechanizm weryfikacji
poprawności. Różnią się natomiast formatem danych treningowych oraz rodzajem
tekstu, który mają wygenerować.

Aby drugi model mógł być trenowany i używany w taki sam sposób jak model
dowodzący, rozbicia muszą zostać zapisane w jednoznacznym formacie tekstowym.
Format ten powinien określać trzy informacje: sekwent, który jest rozbijany, regułę
wnioskowania używaną do rozbicia oraz listę sekwentów-przesłanek otrzymanych po
zastosowaniu tej reguły.

W systemie rozbicie zapisywane jest jako tekst o następującym formacie:
\[
\begin{array}{l}
\texttt{Smash: } S\\
\texttt{With: } R\\
\texttt{To: } S_1,\\
\ldots\\
S_n
\end{array}
\]
gdzie \(S\) jest sekwentem, który chcemy rozbić, \(R\) jest nazwą reguły
wnioskowania, a \(S_1,\ldots,S_n\) są sekwentami-przesłankami otrzymanymi po
rozbiciu.

Pierwsza linia określa problem wejściowy. Druga linia wskazuje regułę, za pomocą
której problem ma zostać rozbity. Kolejne linie zawierają listę podproblemów. Jeżeli
reguła nie wymaga żadnych przesłanek (np reguła \texttt{Assumption}), lista podproblemów jest pusta. 

Przykładowy schemat rozbicia można zapisać następująco:
\[
\begin{array}{l}
\texttt{Smash: } \Gamma \vdash \varphi \to \psi\\
\texttt{With: } \to\texttt{-introduction}\\
\texttt{To: } \Gamma,\varphi \vdash \psi
\end{array}
\]
Taki zapis oznacza, że problem udowodnienia implikacji został zastąpiony problemem
udowodnienia jej prawej strony przy dodatkowym założeniu jej lewej strony.




\subsection{Sprawdzanie poprawności rozbicia}

Ponieważ rozbicia są generowane przez model językowy, każde wygenerowane rozbicie
musi zostać sprawdzone formalnie. Rozbicie składa się z problemu wejściowego \(S\),
reguły \(R\) oraz listy podproblemów \(S_1,\ldots,S_n\). Jest poprawne wtedy, gdy
zastosowanie reguły \(R\) do podproblemów pozwala otrzymać problem wejściowy.

W naszej implementacji reguły wnioskowania mogą być stosowane tylko do sekwentów,
które mają ten sam kontekst bazowy. Dlatego przed sprawdzeniem rozbicia wszystkie
podproblemy sprowadzane są do postaci z pustym kontekstem bazowym. Cały ich
kontekst logiczny zostaje przeniesiony do kontekstu dodatkowego. Jeżeli podproblem
ma kontekst bazowy \(\Gamma_b\), kontekst dodatkowy \(\Gamma_d\) i tezę \(\varphi\),
to zastępujemy go sekwentem
\[
    \emptyset ; \Gamma_b,\Gamma_d \vdash \varphi.
\]
Gdzie średnik oddziela tutaj kontekst bazowy od kontekstu dodatkowego.

Tak samo przekształcany jest problem wejściowy \(S\). Nie jest to potrzebne do
samego zastosowania reguły, ale upraszcza końcowe porównanie: wynik zastosowania
reguły do przekształconych podproblemów porównujemy z problemem wejściowym
zapisanym w tej samej konwencji.

Procedura sprawdzania wygląda więc następująco:


\begin{algorithm}[H]
\caption{Sprawdzanie poprawności rozbicia}
\begin{algorithmic}[1]
\Procedure{SprawdźRozbicie}{$S, R, S_1,\ldots,S_n$}
    \State \(S' \gets\) sekwent \(S\) sprowadzony do postaci z pustym kontekstem bazowym
    \For{\(i = 1,\ldots,n\)}
        \State \(S_i' \gets\) sekwent \(S_i\) sprowadzony do postaci z pustym kontekstem bazowym
        \If{\(S_i\) nie jest sekwentem tautologicznym}
            \State \Return odrzuć rozbicie
        \EndIf
    \EndFor
    \State \(A \gets\) dodatkowe argumenty potrzebne do zastosowania reguły \(R\)
    \State spróbuj zastosować \(R\) do \(S_1',\ldots,S_n'\) z argumentami \(A\)
    \If{zastosowanie reguły zgłasza błąd}
        \State \Return odrzuć rozbicie
    \EndIf
    \State \(T \gets\) wynik zastosowania reguły \(R\)
    \If{\(T \neq S'\)}
        \State \Return odrzuć rozbicie
    \EndIf
    \State \Return zaakceptuj rozbicie
\EndProcedure
\end{algorithmic}
\end{algorithm}



Dodatkowe argumenty podobnie jak przy przetważaniu dowodów liniowych są potrzebne dla tych reguł, których zastosowanie  nie jest
jednoznacznie wyznaczone przez same przesłanki.  W takich przypadkach
brakujące dane są dedukowane z problemu wejściowego lub z problemu wejściowego i przesłanek.

Sprawdzanie tautologiczności podproblemów pełni w systemie rolę pomocniczego
filtra. Nie jest ono potrzebne do samego sprawdzenia, czy dane rozbicie odpowiada
schematowi wybranej reguły wnioskowania. Pozwala jednak od razu odrzucić takie
rozbicia, które prowadzą do podproblemów niedowodliwych, a więc do gałęzi drzewa
poszukiwania, które nie mogłyby zostać domknięte.

Rozwiązanie to jest mało eleganckie z punktu widzenia czysto dowodowego, ponieważ
system sprawdza semantycznie, czy sekwent, dla którego następnie ma znaleźć dowód,
jest tautologiczny. W niniejszej pracy jest to jednak świadome uproszczenie
implementacyjne. System ma charakter prototypowy i eksperymentalny, a rozważane
formuły są na tyle małe, że sprawdzenie wszystkich wartościowań nie stanowi
istotnego kosztu obliczeniowego. Dzięki temu można ograniczyć liczbę oczywiście
bezowocnych prób i skupić się na właściwym celu pracy, czyli generowaniu formalnych
dowodów przez model językowy.

Należy podkreślić, że filtr tautologiczności nie zastępuje dowodu. Zaakceptowany
podproblem nadal musi zostać udowodniony w przyjętej notacji, a każda linia
wygenerowanego dowodu jest następnie sprawdzana przez parser i implementację
reguł wnioskowania. Sprawdzenie tautologiczności pełni więc jedynie funkcję
praktycznego ograniczenia przestrzeni poszukiwań, a nie mechanizmu konstruowania
dowodu.


\subsection{Generowanie rozbić przez model}

Generowanie rozbicia przebiega analogicznie do generowania kolejnych linii dowodu:
model językowy nie tworzy od razu struktury logicznej, lecz dopisuje tekst w ustalonym
formacie. Punktem wyjścia jest prompt zawierający sekwent, który ma zostać rozbity:
\[
\texttt{Smash: } S
\]
Na podstawie tego promptu model najpierw generuje linię wskazującą regułę
wnioskowania, czyli linię zaczynającą się od \texttt{With:}. Następnie generowana jest
linia \texttt{To:}, zawierająca listę podproblemów, na które ma zostać rozbity sekwent
wejściowy.

Podczas generowania podproblemów stosowane jest proste ograniczenie składniowe:
model może dopisywać jedynie symbole logiczne, separatory oraz zmienne występujące
w rozbijanym sekwencie. Dzięki temu zmniejsza się liczba kandydatów, które już na
poziomie zapisu wprowadzają zupełnie nowe zmienne i z tego powodu nie mogą być
sensownymi rozbiciami danego problemu. Nie jest to jednak pełna kontrola poprawności.
Wygenerowany tekst nadal musi zostać sparsowany i sprawdzony przez procedurę
weryfikującą rozbicie.

Model może też wygenerować listę podproblemów w kilku liniach. Przyjmujemy, że opis
rozbicia jest gotowy do sprawdzenia wtedy, gdy ma już co najmniej trzy linie oraz
ostatnia linia nie kończy się przecinkiem. Przecinek na końcu oznacza, że model
sugeruje dalszy ciąg listy podproblemów, więc generowanie powinno być kontynuowane.

Schemat działania można zapisać następująco:

\begin{algorithm}[H]
\caption{Generowanie rozbicia}
\begin{algorithmic}[1]
\Procedure{WygenerujRozbicie}{$S, N$}
    \State \(p_0 \gets \texttt{Smash: } S\)
    \State \(p \gets p_0\)
    \For{\(i = 1,\ldots,N\)}
        \While{\(p\) nie jest gotowe do sprawdzenia}
            \State wygeneruj następną linię tekstu na podstawie \(p\)
            \State dopisz wygenerowaną linię do \(p\)
        \EndWhile

        \If{\(p\) koduje poprawne rozbicie}
            \State \Return \(p\)
        \Else
            \State \(p \gets p_0\)
        \EndIf
    \EndFor
    \State \Return informacja o niepowodzeniu
\EndProcedure
\end{algorithmic}
\end{algorithm}

Jeżeli wygenerowany tekst nie przechodzi weryfikacji, zostaje odrzucony w całości
poza pierwszą linią \(\texttt{Smash:} S\), która określa sekwent do rozbicia. System nie
próbuje naprawiać błędnego kandydata lokalnie, ponieważ błąd może wynikać zarówno
z wyboru niewłaściwej reguły, jak i z niepoprawnej listy podproblemów albo z samego
formatu tekstowego. W takiej sytuacji prostsze i bardziej stabilne jest ponowienie
generacji od tego samego punktu początkowego. Procedura powtarza ten proces do
momentu znalezienia poprawnego rozbicia albo osiągnięcia ustalonego limitu prób.



\subsection{Dowodzenie sekwentów z założeniami}

Dotychczas najprostsza wersja systemu mogła być rozumiana jako narzędzie do
dowodzenia tautologii, czyli sekwentów o pustym kontekście. Wprowadzenie rozbić
problemów na podproblemy zmienia jednak tę sytuację. Nawet jeśli problem wyjściowy
ma postać
\[
    \vdash \varphi,
\]
to jego podproblemy mogą mieć już niepuste konteksty. Przykładowo rozbicie celu
\[
    \vdash \varphi \to \psi
\]
przez regułę wprowadzania implikacji prowadzi do podproblemu
\[
    \varphi \vdash \psi.
\]
Oznacza to, że system musi umieć nie tylko dowodzić tautologii, lecz także generować
dowody sekwentów z założeniami.


Przy dowodzeniu sekwentów z założeniami pojawia się dodatkowa trudność:
należy kontrolować, które formuły należą do kontekstu sekwentu kodowanego przez
daną linię dowodu. Naturalnym uproszczeniem jest przyjęcie, że założenia sekwentu
wejściowego są dostępne przez cały czas trwania dowodu, a więc znajdują się
w kontekście każdej generowanej linii. 

To uproszczenie prowadzi do podziału kontekstu na dwie części. Kontekst bazowy,
oznaczany w implementacji jako \texttt{base\_context}, zawiera formuły traktowane
jako stale dostępne w całym rozważanym dowodzie. Są to założenia
dowodzonego sekwentu, zapisane w nagłówku promptu. Jako \texttt{additional\_context} z kolei oznaczamy kontekst dodatkowy, który zawiera natomiast założenia lokalne,
które mogą zmieniać się pomiędzy liniami dowodu. 

Reguły wnioskowania stosujemy wyłącznie do przesłanek mających ten sam kontekst
bazowy oraz wynik działania reguły ma również ten sam kontekst bazowy co przesłanki. 

Działanie reguł wnioskowania dla przesłanek z kontekstem podzielonym na kontekst bazowy i dodatkowy 
można rozumieć jako zastosowanie tych reguł w klasyczny sposób, czyli
bez rozróżnienia na kontekst bazowy i dodatkowy, a następnie uzupełnieniu kontekstu konkluzji
brakującymi formułami należącymi do kontekstu bazowego. 

Można więc powiedzieć, że schemat działania reguły \(R\) dla sekwentów z rozróżnionym kontekstem bazowym i dodatkowym polega na zastąpieniu schematu

\[
\frac{
    \Gamma_1 \vdash \varphi_1
    \qquad
    \Gamma_2 \vdash \varphi_2
    \qquad
    \cdots
    \qquad
    \Gamma_n \vdash \varphi_n
}{
    \Delta \vdash \psi
}\;(R)
\]

schematem:

\[
\frac{
    \Gamma ; \Gamma_1 \setminus \Gamma \vdash \varphi_1
    \qquad
    \Gamma ; \Gamma_2 \setminus \Gamma \vdash \varphi_2
    \qquad
    \cdots
    \qquad
    \Gamma ; \Gamma_n \setminus \Gamma\vdash \varphi_n
}{
    \Gamma ; \Delta \setminus \Gamma \vdash \psi
}\;(R)
\]

Zakładamy, że kontekst bazowy i dodatkowy każdego sekwentu są rozłączne.

Taka konwencja wymaga jednak uzasadnienia poprawności. Po pierwsze, trzeba
pokazać, że jeśli dany sekwent można udowodnić w zwykłej reprezentacji, czyli bez
wyróżnionego kontekstu bazowego, to można go również udowodnić w reprezentacji,
w której wszystkie sekwenty pośrednie mają jako kontekst bazowy założenia problemu
wejściowego. 

Dowód tej zależności jest bardzo prosty, ponieważ ze zdefiniowania działania reguł dla sekwentów z kontekstem bazowym
wynika wprowst sposób konwersji dowodu z jednej reprezentacji na drugą. Z kolei sekwent w samym korzeniu drzewa dowodu po takiej konwersji ma 
pusty kontekst dodatkowy (bo \(\Gamma ; \Delta \setminus \Gamma \vdash \psi\) przybiera postać \(\Gamma ; \Gamma \setminus \Gamma \vdash \psi\), czyli \(\Gamma ; \emptyset \vdash \psi\)). 

Po drugie, trzeba pokazać zależność odwrotną: użycie kontekstu bazowego nie może
pozwalać na udowodnienie sekwentu, którego nie dałoby się uzyskać w zwykłej
reprezentacji. Jeśli więc istnieje dowód, w którym sekwenty pośrednie korzystają
z niepustego kontekstu bazowego, to powinien istnieć również odpowiadający mu
dowód zapisany bez tego technicznego rozróżnienia, czyli taki, w którym przesłanki mają kontekst bazowy pusty.

Dowód tej zależności będzie polegać na pokazaniu algorytmu, który dowód dla sekwentów z kontekstem bazowym równym kontekstowi bazowemu tezy 
konwertuje na dowód z sekwentami o pustym kontekście bazowym.

Dla uproszczenia rozważań przyjmujemy konwencję jest w której istnieje ukryta linia o indeksie \texttt{0}. Nie jest ona
wypisywana w tekście dowodu, a numeracja jawnych linii zaczyna się od \texttt{1}.
Linia \texttt{0} koduje sekwent, którego kontekstem jest przyjęty kontekst bazowy,
a tezą jest \(\top\). Dzięki temu reguły \texttt{context} oraz \texttt{contextWeak}
mogą odwoływać się do całego kontekstu bazowego bez konieczności dopisywania
osobnej linii dowodu której jedyną funkcją jest posiadanie kontekstu równego koqntekstowi bazowemu. 
Jeżeli kontwekst bazowy jest pusty, linia \texttt{0} koduje sekwent \(\emptyset \vdash \top\).

Idea algorytmu bazuje na następującej obserwacji: jeżeli poprawny jest schemat 


\[
\frac{
    \Gamma ; \Gamma_1 ' \vdash \varphi_1
    \qquad
    \Gamma ; \Gamma_2 '\vdash \varphi_2
    \qquad
    \cdots
    \qquad
    \Gamma ; \Gamma_n ' \vdash \varphi_n
}{
    \Gamma ; \Delta ' \vdash \psi
}\;(R)
\]

To z definicji działania reguł dla sekwentów z kontekstem bazowym wynika, że istnieją takie zbiory formuł \(\Gamma_1,\ldots,\Gamma_n\) oraz \(\Delta\), że 
\(\Gamma_i ' = \Gamma_i \setminus \Gamma\) dla \(i=1,\ldots,n\) oraz \(\Delta ' = \Delta \setminus \Gamma\) oraz schemat


\[
\frac{
    \Gamma_1  \vdash \varphi_1
    \qquad
    \Gamma_2  \vdash \varphi_2
    \qquad
    \cdots
    \qquad
     \Gamma_n  \vdash \varphi_n
}{
    \Delta \vdash \psi
}\;(R)
\]

Jednak ze schematu każdej reguły wnioskowania w naszym systemie, która ma niepustą listę przesłanek, 
wynika, że jeżeli dla przesłanek \(S_1,\ldots,S_n\) można zastosować regułę \(R\) 
można ją również zastosować do tych samych przesłanek, ale z rozszerzonymi o jakieś formuły kontekstami, a co za tym idzie w rozważanym przypadku


\[
\frac{
    \Gamma_1 \cup \Gamma \vdash \varphi_1
    \qquad
    \Gamma_2 \cup \Gamma \vdash \varphi_2
    \qquad
    \cdots
    \qquad
    \Gamma_n \cup \Gamma \vdash \varphi_n
}{
    \Gamma ' \vdash \psi
}\;(R)
\]

dla jakiegoś \(\Gamma'\) zawierającego się w \(\Gamma \cup \Delta\).




Algorytm usuwania kontekstu bazowego działa następująco. Niech
\[
    \Gamma = \{\gamma_1,\ldots,\gamma_m\}
\]
będzie kontekstem bazowym dowodu wejściowego. Na początku nowego dowodu
wprowadzamy kolejno wszystkie formuły z \(\Gamma\) jako jawne linie dowodu.
Robimy to przy użyciu reguły \texttt{contextWeak}. Pierwsza z tych linii odwołuje
się do ukrytej linii o indeksie \texttt{0}, a każda kolejna do poprzednio uzyskanej linii, dzięki
czemu po \(m\) krokach otrzymujemy linię, której kontekst dodatkowy jest równy całemu
dawnemu kontekstowi bazowemu (można to zrobić bo linii o indeksie \texttt{0} używamy tylko w regule \texttt{context} i 
\texttt{contextWeak}, dla których nie jest istotna konkluzja przesłanki).

Następnie algorytm przechodzi przez kolejne linie oryginalnego dowodu, ale nie
kopiuje ich dosłownie. Dla każdej starej linii odczytywana jest reguła, która została
w niej użyta, oraz indeksy wcześniejszych linii, do których ta reguła się odwołuje.
Ponieważ w nowym dowodzie pojawiają się dodatkowe linie, stare indeksy nie mogą
zostać użyte bezpośrednio. Algorytm przechowuje więc mapę indeksów, która każdej
linii starego dowodu przypisuje linię odpowiadającą jej w nowym dowodzie.

Dla rozważanej starej linii algorytm bierze indeksy jej przesłanek, zastępuje je
indeksami wskazanymi przez tę mapę, a następnie stosuje tę samą regułę do
odpowiadających im linii nowego dowodu. Otrzymana w ten sposób linia nie jest
kopią starej linii, lecz jej odpowiednikiem skonstruowanym już w dowodzie bez
kontekstu bazowego. Jeżeli stara linia odwoływała się do ukrytej linii o indeksie \texttt{0},
to odwołanie to zostaje zastąpione odwołaniem do linii o indeksie \texttt{m} w nowym dowodzie, czyli do linii, której kontekst dodatkowy
 jest równy całemu dawnemu kontekstowi bazowemu (można to zrobić bo linii o indeksie \texttt{0} używamy tylko w regule \texttt{context} i 
\texttt{contextWeak}, dla których nie jest istotna konkluzja przesłanki).

Po skonstruowaniu takiego odpowiednika algorytm sprawdza, czy jego kontekst
dodatkowy zawiera wszystkie formuły z dawnego kontekstu bazowego. Jeżeli którejś
z nich brakuje, dodawane są kolejne linie otrzymane przez zastosowanie reguły
\texttt{weakening}. Dopiero ostatnia linia uzyskana po tych ewentualnych osłabieniach
jest traktowana jako właściwy odpowiednik danej linii starego dowodu i jej numer
zostaje zapisany w mapie indeksów.

Dla skrócenia otrzymanego dowodu, jeżeli przetważana stara linia była linią z regułą \texttt{assumption},
to w nowym dowodzie stosowana jest reguła \texttt{contextWeak} odwołująca się do linii o indeksie \texttt{m}.
  
Schematycznie procedurę można zapisać następująco:

\begin{algorithm}[H]
\caption{Usuwanie kontekstu bazowego}
\label{alg:remove-base-context}
\begin{algorithmic}[1]
\Procedure{UsuńKontekstBazowy}{$D$}
    \State $\Gamma=\{\gamma_1,\ldots,\gamma_m\} \gets$ kontekst bazowy dowodu $D$
    \State $D' \gets$ pusty dowód
    \State $f(0) \gets 0$

    \For{$i = 1,\ldots,m$}
        \State dopisz do $D'$ linię wyprowadzającą $\gamma_i$ przypomocy \texttt{contextWeak, i-1}
    \EndFor

    \State $f(0) \gets m$

    \For{każdej jawnej linii $l_i$ dowodu $D$}
        \State odczytaj z $l_i$ regułę $R$, formułę $\varphi_i$ oraz indeksy $j_1,\ldots,j_n$
        \State zamień indeksy przesłanek na $f(j_1),\ldots,f(j_n)$

        \If{$R=\texttt{assumption}$}
            \State dopisz do $D'$ wynik zastosowania \texttt{contextWeak} do linii o indeksie $f(0)$
        \EndIf
     
        \If{$R \neq \texttt{assumption}$}
            \State dopisz do $D'$ wynik zastosowania $R$ do linii $f(j_1),\ldots,f(j_n)$

        \For{$j \in 1,\ldots,m$}
            \If{$\gamma_j$ nie należy do kontekstu ostatniej linii $D'$}
                \State $n \gets$ liczba linii $D'$
                \State \texttt{l} = \texttt{n+1. }$\phi_i$ \texttt{ weakening j, n}
                \State dopisz \texttt{l} do $D'$ 
            \EndIf
        \EndFor
        \EndIf
          
        

        \State $f(i) \gets$ numer ostatniej linii dowodu $D'$
    \EndFor

    \State \Return $D'$
\EndProcedure
\end{algorithmic}
\end{algorithm}





\subsection{Drzewo poszukiwań}

W tej części pracy opisujemy główny algorytm strategii trzeciej, oraz strukturę danych, na której jest ona oparta.

Aby połączyć dekompozycję celu z generowaniem dowodów częściowych, system
buduje drzewo poszukiwania. Korzeniem tego drzewa jest sekwent, który chcemy
udowodnić. Następnie system może wygenerować kilka różnych rozbić tego sekwentu
na podproblemy. Każde takie rozbicie stanowi osobną możliwość kontynuacji
poszukiwania dowodu. Dzięki temu system nie musi od razu decydować się na jedną
postać rozumowania, lecz może przechowywać równolegle kilka alternatywnych
struktur dowodu.

W drzewie występują dwa typy wierzchołków. Pierwszy typ reprezentuje problemy
dowodowe, czyli sekwenty, które należy udowodnić. Będziemy oznaczać je jako
wierzchołki typu \texttt{LP}. Drugi typ reprezentuje konkretne rozbicia problemu na
podproblemy. Będziemy oznaczać je jako wierzchołki typu \texttt{SP}.

Dzieci wierzchołka typu \texttt{LP} są wierzchołkami typu \texttt{SP}, kodującymi
możliwe rozbicia sekwentu zapisanego w ich rodzicu. Z kolei dzieci wierzchołka typu
\texttt{SP} są wierzchołkami typu \texttt{LP}, które kodują sekwenty będące
przesłankami rozbicia zapisanego w ich rodzicu. Struktura drzewa jest więc
naprzemienna: dziećmi wierzchołków problemów dowodowych są wierzchołki rozbić, 
a dziećmi wierzchołków rozbić są wierzchołki problemów dowodowych.

Intuicyjnie działanie algorytmu polega na wielokrotnym wybieraniu nierozwiązanych
fragmentów drzewa i próbie ich dowiedzenia. 

Jeżeli próbujemy udowodnić wierzchołek typu \texttt{LP}, który jest liściem, system najpierw próbuje znaleźć
jego bezpośredni dowód za pomocą strategii drugiej. Jeśli to się uda, dowód
zostaje zapisany w wierzchołku, a wartość wierzchołka zmienia się na prawdziwą.
Jeżeli bezpośredni dowód nie zostanie znaleziony, system generuje możliwe rozbicia
tego sekwentu i dodaje je jako dzieci wierzchołka. Następnie jako dzieci każdego
z tych dzieci dodawane są wierzchołki typu \texttt{LP} kodujące przesłanki
rozbicia zapisanego w ich rodzicu. Jeżeli któryś z wierzchołków rozbicia nie ma
żadnych przesłanek, zostaje od razu oznaczony jako udowodniony.

Wierzchołek typu \texttt{SP} reprezentuje jedno konkretne rozbicie. Jego dzieciami są
sekwenty-przesłanki tego rozbicia. Aby udowodnić rozbicie, trzeba udowodnić
wszystkie te sekwenty. Dlatego dla wierzchołka typu \texttt{SP} algorytm próbuje
rekurencyjnie dowieść wszystkich jego dzieci. Jeśli choć jedno z nich nie zostanie
udowodnione w danej próbie, całe rozbicie nie zostaje uznane za udowodnione.

W przypadku wewnętrznego wierzchołka typu \texttt{LP} sytuacja jest inna: wystarczy,
że powiedzie się próba udowodnienia jednego z jego rozbić. Dlatego algorytm wybiera
jedno z jego dzieci typu \texttt{SP} i próbuje je udowodnić. Wybór ten odbywa się
z wagami wprost proporcjonalnymi do częstości, z jaką dane rozbicie zostało
wygenerowane przez model. Jeżeli wybrane rozbicie zostanie udowodnione, to
udowodniony zostaje również sekwent reprezentowany przez jego rodzica.

Po każdej udanej próbie dowodu wierzchołka informacja o tym które wierzchołki zostały udowodnione
jest propagowana w górę drzewa. Jeżeli wierzchołek typu \texttt{SP} zostaje udowodniony, 
jego rodzic typu \texttt{LP} jest oznaczany jako udowodniony, a informacja o tym jest przekazywana 
dalej do jego rodzica, aż do korzenia drzewa.
Jeżeli zaś udowodniony zostaje wierzchołek typu \texttt{LP}, to dla jego rodzica typu \texttt{SP} 
sprawdzamy, czy wszystkie jego dzieci są już udowodnione. Jeżeli tak, to również ten wierzchołek typu \texttt{SP}
jest oznaczany jako udowodniony, a informacja o tym jest propagowana dalej dla kolejnych przodków, aż do korzenia.

Opisany mechanizm można wygodnie podzielić na cztery procedury. Pierwsza z nich
odpowiada za obsługę liścia typu \texttt{LP}, czyli sekwentu, dla którego nie znamy
jeszcze ani bezpośredniego dowodu, ani rozbicia na podproblemy. Druga procedura
odpowiada za samo wygenerowanie rozbić takiego liścia i dołączenie ich do drzewa.
Trzecia procedura opisuje ogólne rekurencyjne przeszukiwanie drzewa, korzystające
z dwóch poprzednich procedur jako podprogramów. Czwarta procedura zaś uruchamia w pętli 
procedurę trzecią poszukując dowodu w drzewie po kolejnych aktualizacjach jego struktury.

Najpierw rozważmy obsługę liścia typu \texttt{LP}. Taki wierzchołek reprezentuje
pojedynczy sekwent, który system próbuje udowodnić. Naturalnie najpierw sprawdzamy,
czy da się znaleźć jego dowód bez dalszego rozbijania problemu. W tym celu
uruchamiana jest strategia druga szukania dowodu.
Jeżeli zakończy się ona sukcesem, znaleziony dowód zostaje zapisany w wierzchołku,
a sam wierzchołek zostaje oznaczony jako udowodniony. Jeżeli natomiast dowodu
bezpośredniego nie uda się znaleźć w przyjętych limitach, wierzchołek zostaje
rozbity na podproblemy.

\begin{algorithm}[H]
\caption{Obsługa liścia typu \texttt{LP}}
\label{alg:run-leaf}
\begin{algorithmic}[1]
\Procedure{ObsluzLisc}{$X$}
    \State spróbuj znaleźć bezpośredni dowód sekwentu kodowanego przez $X$

    \If{dowód został znaleziony}
        \State zapisz znaleziony dowód w $X$
        \State oznacz $X$ jako udowodniony
        \State zaktualizuj wartości przodków $X$
        \State \Return True
    \Else
        \State \Call{RozbijLisc}{$X$}
        \State \Return False
    \EndIf
\EndProcedure
\end{algorithmic}
\end{algorithm}

Jeżeli nie uda się znaleźć bezpośredniego dowodu sekwentu zapisanego w liściu,
procedura obsługi liścia deleguje rozbicie tego wierzchołka do kolejnej procedury.
Procedura rozbijania nie wybiera zwykle pojedynczego rozbicia, lecz wykonuje
określoną liczbę prób wygenerowania rozbicia przez model. Liczbę tych prób
oznaczamy przez \(N_{\mathrm{roz}}\). Każde poprawne rozbicie jest zapamiętywane,
a jeśli to samo rozbicie pojawi się wielokrotnie (co może zdarzać się dość często, bo często jest
bardzo mało naturalnych rozbić), zwiększana jest jego waga używana w kolejnej procedurze do losowania obsługiwanego wierzchołka. Dzięki
temu powtarzające się propozycje modelu nie są traktowane jako zbędne duplikaty,
lecz jako informacja o tym, które rozbicia model uznaje za bardziej naturalne.

Po zakończeniu generowania poprawne rozbicia są dodawane jako dzieci rozważanego
wierzchołka typu \texttt{LP}. Następnie dla każdego wierzchołka typu \texttt{SP}
dodawane są jego dzieci typu \texttt{LP}, kodujące sekwenty-przesłanki danego
rozbicia. Jeżeli dane rozbicie nie ma żadnych przesłanek, odpowiadający mu
wierzchołek typu \texttt{SP} zostaje od razu oznaczony jako udowodniony.


\begin{algorithm}[H]
\caption{Rozbijanie liścia typu \texttt{LP}}
\label{alg:smash-leaf}
\begin{algorithmic}[1]
\Procedure{RozbijLisc}{$X, N_{\mathrm{roz}}$}
    \State utwórz prompt opisujący sekwent kodowany przez $X$
    \State $W \gets$ pusty słownik wag rozbić

    \For{$i = 1,\ldots,N_{\mathrm{roz}}$}
        \State wygeneruj kandydackie rozbicie $S$
        \If{$S$ jest poprawnym rozbiciem}
            \State zwiększ wagę rozbicia $S$ w słowniku $W$
        \EndIf
    \EndFor


    \For{każde rozbicie $S$ zapisane w $W$}
        \State dodaj do $X$ dziecko typu \texttt{SP} kodujące rozbicie $S$
        \State przypisz temu dziecku wagę proporcjonalną do częstości $S$
    \EndFor

    \For{każde dziecko $Y$ wierzchołka $X$}
        \For{każda przesłanka $P$ rozbicia zapisanego w $Y$}
            \State dodaj do $Y$ dziecko typu \texttt{LP} kodujące sekwent $P$
        \EndFor

        \If{$Y$ nie ma dzieci}
            \State oznacz $Y$ jako udowodniony
            \State zaktualizuj wartości przodków $Y$
        \EndIf
    \EndFor
\EndProcedure
\end{algorithmic}
\end{algorithm}


Trzecia procedura opisuje właściwe przeszukiwanie drzewa. Jeżeli aktualny
wierzchołek jest już oznaczony jako udowodniony, algorytm od razu kończy pracę
sukcesem. Jeżeli jest liściem typu \texttt{LP}, uruchamiana jest opisana wyżej
procedura obsługi liścia. Jeżeli jest wewnętrznym wierzchołkiem typu \texttt{LP},
algorytm wybiera jedno z jego dzieci typu \texttt{SP} zgodnie z wagami rozbić i próbuje
udowodnić właśnie tę gałąź. Jeżeli natomiast aktualny wierzchołek jest typu
\texttt{SP}, to algorytm musi udowodnić wszystkie jego dzieci, ponieważ dane rozbicie
jest skuteczne dopiero wtedy, gdy rozwiązane są wszystkie jego podproblemy.

\begin{algorithm}[H]
\caption{Przeszukiwanie drzewa dowodu}
\label{alg:try-node}
\begin{algorithmic}[1]
\Procedure{Dowiedz}{$X$}
    \If{$X$ jest już oznaczony jako udowodniony}
        \State \Return True
    \EndIf

    \If{$X$ jest liściem typu \texttt{LP}}
        \State \Return \Call{ObsluzLisc}{$X$}
    \EndIf

    \If{$X$ jest wewnętrznym wierzchołkiem typu \texttt{LP}}
        \State wybierz jedno z dzieci $X$, oznaczmy je przez $Y$
        \State wybór wykonaj z prawdopodobieństwami zapisanymi w dzieciach $X$
        \State \Return \Call{Dowiedz}{$Y$}
    \EndIf

    \If{$X$ jest wierzchołkiem typu \texttt{SP}}
        \For{każde dziecko $Y$ wierzchołka $X$}
            \State $b \gets$ \Call{Dowiedz}{$Y$}
            \If{$b = False$}
                \State \Return False
            \EndIf
        \EndFor

        \State oznacz $X$ jako udowodniony
        \State zaktualizuj wartości przodków $X$
        \State \Return True
    \EndIf
\EndProcedure
\end{algorithmic}
\end{algorithm}

Główna procedura poszukiwania dowodu uruchamia procedurę \textsc{Dowiedz}
w pętli. Pojedyncze wywołanie nie musi od razu znaleźć całego dowodu, a może jedynie rozbudować drzewo lub oznaczyć część
wierzchołków jako udowodnione, dlatego kolejne wywołania korzystają z coraz bardziej
rozwiniętej struktury. Działanie kończy się sukcesem, jeżeli korzeń drzewa zostanie
oznaczony jako udowodniony. Wtedy system odtwarza z drzewa pełny dowód tekstowy.
Jeżeli nie uda się tego zrobić po \(N_{\max}\) próbach, procedura kończy się
niepowodzeniem.

\begin{algorithm}[H]
\caption{Główna pętla poszukiwania dowodu}
\label{alg:main-proof-loop}
\begin{algorithmic}[1]
\Procedure{SzukajDowodu}{$G, N_{\max}$}
    \State utwórz korzeń $K$ typu \texttt{LP} kodujący cel $G$

    \For{$i = 1,\ldots,N_{\max}$}
        \If{$K$ jest oznaczony jako udowodniony}
            \State \Return \Call{OdtworzDowod}{$K$}
        \EndIf

        \State $b \gets$ \Call{Dowiedz}{$K$}

        \If{$b = True$ albo $K$ jest oznaczony jako udowodniony}
            \State \Return \Call{OdtworzDowod}{$K$}
        \EndIf
    \EndFor

    \State \Return False
\EndProcedure
\end{algorithmic}
\end{algorithm}


\subsection{Sklejanie dowodów}

Samo znalezienie rozwiązania w drzewie poszukiwania nie oznacza jeszcze, że mamy
gotowy tekstowy dowód pierwotnej tezy. Drzewo zawiera informację o tym, które
podproblemy zostały udowodnione oraz przez jakie rozbicia można je połączyć, ale
sam dowód jest rozproszony po różnych wierzchołkach. 

Procedura sklejania działa rekurencyjnie od liści drzewa ku korzeniowi. Jeżeli
rozwiązany wierzchołek typu \texttt{LP} jest liściem, to zawiera bezpośrednio
wygenerowany dowód danego sekwentu. Taki dowód może mieć niepusty kontekst
bazowy, dlatego najpierw stosowana jest opisana wcześniej procedura usuwania
kontekstu bazowego. W efekcie otrzymujemy zwykły fragment dowodu, który można
wkleić do większego dowodu.

Jeżeli natomiast wierzchołek typu \texttt{LP} nie jest liściem, to znaczy, że został
rozwiązany przez jedno ze swoich rozbić. Wtedy procedura przechodzi do tego dziecka
typu \texttt{SP}, które zostało oznaczone jako udowodnione. Dla wierzchołka typu
\texttt{SP} najpierw rekurencyjnie odtwarzane są dowody wszystkich jego dzieci, czyli
poddowody przesłanek rozbicia. Następnie te poddowody są łączone w jeden ciąg
linii, a na końcu dopisywana jest linia stosująca regułę zapisaną w rozbiciu.

Oczywiście łączenie dowodów wymaga odpowiedniej numeracji linii. Numeracja w każdym poddowodzie
rozpoczynana jest od \texttt{1}, dlatego
przed połączeniem kolejnych fragmentów wszystkie indeksy w danym poddowodzie są
przesuwane o liczbę linii występujących we wcześniejszych fragmentach. Dotyczy to
zarówno numerów samych linii, jak i odwołań do wcześniejszych linii zapisanych przy
regułach.

Jeżeli mamy poddowody
\[
D_1,\ldots,D_n,
\]
ich ostatnie linie kodują sekwenty będące
przesłankami reguły, przez którą dane podproblemy zostały połączone w drzewie
poszukiwania. Numery tych linii (po przenumerowaniu) mogą więc zostać użyte jako odwołania w końcowej
linii sklejającej dowód.


Kolejność tych numerów musi odpowiadać kolejności przesłanek danej reguły. W implementacji
zapewnia to struktura drzewa poszukiwania: dzieci węzła reprezentującego rozbicie
są przechowywane w tej samej kolejności, w jakiej odpowiadają przesłankom reguły
użytej do rozbicia. Dzięki temu, po sklejeniu poddowodów, można odczytać numery
ich końcowych linii jako uporządkowaną listę argumentów końcowej reguły.

Dla większości reguł wystarczy więc połączyć przenumerowane poddowody, a następnie
dopisać jedną linię postaci
\[
k_n+1\texttt{. }\text{konkluzja rozbicia}
\qquad
R,\; k_1, \ldots, k_n,
\]
gdzie \(k_1,\ldots,k_n\) są numerami końcowych linii odpowiednich poddowodów. Separatory występujące
w uzasadnieniu linii są przy tym generowane zgodnie z formatem tekstowym danej
reguły. W omawianych tu przypadkach są to przecinki oddzielające kolejne indeksy.

Niektóre reguły wymagają jednak dodatkowej ostrożności. W implementacji wyróżnione
są reguły, których zapis tekstowy odwołuje się nie tylko do linii będących
bezpośrednimi przesłankami reguły, ale także do dodatkowych linii pomocniczych.


Dotyczy to reguł:
\[
\texttt{weakening}, \qquad
\to\texttt{-introduction}, \qquad
\neg\texttt{-introduction}, \qquad
\texttt{RAA}.
\]
Dla takich reguł nie wystarcza wskazanie samych końcowych linii poddowodów.
Trzeba jeszcze dopisać odpowiednią linię pomocniczą, do której końcowa reguła będzie
mogła się odwołać. W tym celu każdą z tych linii rozpatrujemy osobno.

Przykładowo dla reguły wyprowadzania implikacji końcowa linia ma postać
\[
\texttt{i. }
    \varphi \to \psi
    \qquad
    \to\texttt{-introduction j–k}, 
\]
Indeks \texttt{k} wskazuje linię przesłanki, z kolei indeks \texttt{j} wskazuje linię,
z której odczytujemy argument \(\varphi\) którego nie jesteśmy w stanie jednoznacznie zidentyfikować
na podstawie samej przesłanki. W tym celu do złączonych dowodów przesłanek dodajemy linię, która przy pomocy reguły \texttt{assumption}
wyprowadza argument \(\varphi\). Indeks tej linii jest następnie używany jako indeks \texttt{j} w końcowej linii.
Podobne schematy konstruujemy dla reguł \texttt{weakening}, \texttt{negation-introduction} oraz \texttt{RAA}.


Ogólny schemat sklejania poddowodów odpowiadających jednemu rozbiciu można
zapisać następująco:

\begin{algorithm}[H]
\caption{Sklejanie poddowodów jednego rozbicia}
\label{alg:glue-up}
\begin{algorithmic}[1]
\Procedure{Sklej}{$D_1,\ldots,D_n, S$}
    \State \(R \gets\) reguła zapisana w rozbiciu \(S\)
    \State \(P \gets\) problem rozbijany przez \(S\)


    \State \(B \gets\) pusty tekst dowodu
    \State \(k \gets 0\)

    \For{\(i = 1,\ldots,n\)}
        \If{\(D_i\) ma niepusty kontekst bazowy}
            \State \(D_i \gets\) wynik usunięcia kontekstu bazowego z \(D_i\)
        \EndIf

        \State przesuń wszystkie indeksy w \(D_i\) o \(k\)
        \State dopisz \(D_i\) do \(B\)
        \State \(k \gets\) liczba linii aktualnego dowodu \(B\)
    \EndFor

    \If{\(R\) nie wymaga dodatkowej linii pomocniczej}
        \State dopisz do \(B\) linię stosującą \(R\) do końcowych linii poddowodów
    \Else
        \State dopisz odpowiednią linię pomocniczą
        \State dopisz do \(B\) linię stosującą \(R\) dla linii pomocniczej i końcowych linii poddowodów
    \EndIf
    \State \Return \(B\)
\EndProcedure
\end{algorithmic}
\end{algorithm}

Całe odtwarzanie dowodu z drzewa można opisać osobną procedurą. Jeżeli wierzchołek
typu \texttt{LP} jest rozwiązanym liściem, zwracamy zapisany w nim dowód po usunięciu
kontekstu bazowego. Jeżeli wierzchołek typu \texttt{LP} został rozwiązany przez jedno
z rozbić, przechodzimy do tego dziecka typu \texttt{SP} i w nim odtwarzamy dowód. Jeżeli wierzchołek jest typu
\texttt{SP}, najpierw odtwarzamy dowody wszystkich jego dzieci, a potem sklejamy je
przy użyciu rozbicia zapisanego w tym wierzchołku.

\begin{algorithm}[H]
\caption{Odtwarzanie dowodu z drzewa}
\label{alg:compose-proof}
\begin{algorithmic}[1]
\Procedure{OdtworzDowod}{$X$}
    \If{\(X\) nie jest oznaczony jako udowodniony}
        \State zgłoś błąd
    \EndIf

    \If{\(X\) jest liściem typu \texttt{LP}}
        \State \Return dowód zapisany w \(X\) po usunięciu kontekstu bazowego
    \EndIf

    \If{\(X\) jest wewnętrznym wierzchołkiem typu \texttt{LP}}
        \State \(Y\) \(\gets\) udowodnione dziecko \(X\)
        \State \Return \Call{OdtworzDowod}{$Y$}
    \EndIf

    \If{\(X\) jest wierzchołkiem typu \texttt{SP}}
        \State \(L \gets\) pusta lista dowodów

        \For{każde dziecko \(Y\) wierzchołka \(X\)}
            \State dopisz \Call{OdtworzDowod}{$Y$} do listy \(L\)
        \EndFor

        \State \Return \Call{Sklej}{$L, X$}
    \EndIf
\EndProcedure
\end{algorithmic}
\end{algorithm}

W ten sposób procedura sklejania odwraca proces dekompozycji celu. Rozbijanie
zastępuje jeden problem listą podproblemów, natomiast sklejanie bierze dowody tych
podproblemów i przy użyciu reguły zapisanej w rozbiciu tworzy dowód problemu
wyjściowego. Dzięki temu wynik działania całej strategii nie jest jedynie informacją,
że teza została znaleziona w drzewie poszukiwania, ale konkretnym dowodem zapisanym
w tej samej notacji tekstowej, która była używana przez parser, weryfikator i model
językowy.

\chapter{Generowanie korpusu}

\section{Motywacja}

Wykorzystanie strategii opisanych w poprzednim rozdziale wymaga dwóch odpowiednio
wytrenowanych modeli językowych. Pierwszy z nich ma generować poprawne
dowody z założeniami, w tym również dowody z pustą listą założeń, drugi zaś
rozbicia problemów dowodowych na podproblemy. Do ich wytrenowania potrzebne są
więc duże korpusy danych zawierające przykłady takich dowodów i rozbić zapisane
w przyjętej notacji tekstowej.

Naturalnym punktem wyjścia byłoby znalezienie gotowej bazy tautologii wraz z ich
dowodami, a następnie przekształcenie tych dowodów do formatu używanego w
niniejszej pracy. W praktyce takie podejście okazało się niewystarczające. Większość
dostępnych zbiorów tautologii ma charakter dydaktyczny: są to zwykle zestawy
ćwiczeń dla osób uczących się logiki, zawierające co najwyżej kilkadziesiąt
przykładów. Taka liczba dowodów jest zbyt mała, aby mogła stanowić samodzielny
korpus treningowy dla modelu językowego. Z tego powodu konieczne było
wygenerowanie korpusu danych w sposób algorytmiczny.

Ogólny schemat generowania korpusu był następujący. Najpierw przygotowano
dowody tautologii pochodzących ze skryptu do kursu \emph{Logika dla informatyków}
na Uniwersytecie Wrocławskim. Następnie dowody te połączono w jeden duży ciąg
poprawnych linii dowodowych, który dalej rozszerzano przez algorytmiczne dopisywanie
kolejnych losowych, ale formalnie poprawnych linii.

Z tak otrzymanego dużego ciągu wyłuskiwano następnie dowody tautologii, czyli
dowody linii kodujących sekwenty o pustym kontekście. Otrzymane dowody były
dalej przekształcane w przykłady treningowe używane w różnych częściach systemu:
dowody z założeniami oraz rozbicia problemów dowodowych na podproblemy. Na
końcu korpusy były iteracyjnie formatowane, filtrowane i uzupełniane na podstawie
wyników treningu, tak aby model możliwie dobrze uczył się generowania potrzebnych
fragmentów dowodów. Szczegółowe omówienie poszczególnych etapów znajduje się w
kolejnych sekcjach tego rozdziału.

\section{Formuły źródłowe i początkowe dowody}

Pierwszy etap generowania korpusu polegał na przygotowaniu początkowego zbioru
poprawnych dowodów tautologii. Zbiór ten miał pełnić rolę punktu wyjścia dla
późniejszego algorytmicznego rozszerzania danych.

Na początku przygotowano listę formuł rachunku zdań pochodzących ze skryptu do
kursu \emph{Logika dla informatyków} na Uniwersytecie Wrocławskim. Część z nich była
tautologiami, część formułami spełnialnymi, lecz nietautologicznymi, a część formułami
sprzecznymi. Przy pomocy prostego algorytmu sprawdzania tautologiczności przez
podstawianie wartości logicznych wybrano formuły tautologiczne. Do listy dodano
również negacje formuł sprzecznych, ponieważ one także są tautologiami.

Kolejnym etapem było przygotowanie dowodów otrzymanych tautologii. 
W tym celu wykorzystano zewnętrzny model językowy, któremu przekazano opis używanej
notacji. Następnie ze względu na możliwość halucynacji modelu generującego dowody 
każdy dowód został formalnie sprawdzony przy pomocy weryfikatora omawianego w poprzednim rozdziale.

Po weryfikacji początkowe dowody zostały połączone w jeden duży ciąg linii
dowodowych, który będziemy w dalszej części nazywać "dużym dowodem". Ciąg ten posłużył jako punkt wyjścia dla procedury dalszego
rozszerzania korpusu.
\section{Losowe rozszerzanie dużego dowodu}

Kolejnym etapem generowania korpusu było dopisywanie nowych linii do dużego
dowodu. Ogólna idea tej procedury jest prosta: najpierw losowana jest jedna z reguł
wnioskowania, a następnie wybierane są wcześniejsze linie, które mogą zostać użyte
jako argumenty tej reguły. Jeżeli dana reguła wymaga dodatkowego argumentu, który
nie wynika jednoznacznie z wybranych przesłanek, argument ten jest generowany
osobno. Przykładowo dla linii postaci
\[
\texttt{i. } \varphi \texttt{    assumption}
\]
konieczne jest wygenerowanie formuły \(\varphi\).

Najprostszą metodą tej procedury polegałoby na losowym wybieraniu potencjalnych
argumentów reguły, konstruowaniu kandydackiej linii, a następnie sprawdzaniu, czy
dowód rozszerzony o tę linię poprawnie się parsuje. Jeśli parsowanie kończyło się
sukcesem, linia była dopisywana do dużego dowodu; w przeciwnym razie kandydatura
była odrzucana.

W praktyce takie podejście opierające się na tej ideii okazało się jednak zbyt wolne. Dla części reguł już przy
stosunkowo niewielkim rozmiarze dużego dowodu prawdopodobieństwo wylosowania
pasujących argumentów było małe. W efekcie znalezienie jednej poprawnej kandydackiej
linii mogło wymagać bardzo dużej liczby prób, co znacząco spowalniało cały proces
generowania danych.

Możliwym rozwiązaniem byłoby wprowadzić maksymalną liczbę prób dla danej reguły i po jej
przekroczeniu porzucać próbę dopisania linii, przechodząc do losowania następnej
reguły. Takie rozwiązanie prowadziłoby jednak do niekorzystnego rozkładu przykładów
w korpusie. Reguły trudniejsze do zastosowania, wymagające spełnienia większej
liczby warunków, pojawiałyby się znacznie rzadziej, mimo że często są istotne w
rzeczywistych dowodach. Jednocześnie nadmiernie reprezentowane byłyby reguły
bardzo proste, takie jak \(\texttt{assumption}\), czy szczególnie trywialne w procesie generacji ale
rzadko potrzebne w praktyce, takie jak \(\top\texttt{-introduction}\).

Prowadziłoby to do sytuacji szczególnie niekorzystnej z punktu widzenia treningu
modelu: reguły trudne do nauczenia i ważne dla konstrukcji dowodów byłyby słabo
reprezentowane w danych, natomiast reguły proste lub mało informacyjne dominowałyby
w korpusie. Z tego powodu konieczne było zoptymalizowanie samego procesu doboru
argumentów dla losowanych reguł.

Naturalnym kierunkiem optymalizacji jest zawężenie przestrzeni przeszukiwań, by 
zwiększyć prawdopodobieństwo wylosowania pasujących argumentów. Np dla reguły 


\[
\begin{array}{l}
\texttt{j. } \phi \quad \dots \\
\quad \dots \\
\texttt{k. } \phi \to \psi \quad \dots \\
\quad \dots \\
\texttt{i. } \psi \qquad \texttt{→-elimination, k, j}
\end{array}
\]

której zapis w notacji drzewiastej ma postać

$$
\frac{\Gamma_1 \vdash \varphi \to \psi \qquad \Gamma_2 \vdash \varphi}
     {\Gamma_1,\Gamma_2 \vdash \psi}\;(\to E)
$$

chcielibyśmy na początku ograniczyć się do losowania linii, których konkluzja jest implikacją,
a następnie do linii, której konkluzja jest lewą stroną tej implikacji.

W tym celu linie dużego dowodu są grupowane w kilku słownikach. Pierwszy
słownik grupuje indeksy linii według typu tezy sekwentu. Kluczami są tu rodzaje formuł, np.
koniunkcja, alternatywa, negacja, szczególny przypadek podwójnej negacji czy formuła atomowa,
z kolei wartościami są indeksy linii danego typu.
Drugi słownik grupuje indeksy według dokładnej postaci tezy. Kluczami są tu konkretne formuły,
a wartościami indeksy linii, których tezą jest dana formuła. Trzeci słownik grupuje indeksy linii według
założeń występujących w kontekście. Kluczami są tu formuły, a wartościami indeksy linii, których kontekst zawiera daną formułę.
Ostatnią pomocniczą strukturą jest zestaw indeksów linii, które mogą być użyte jako przesłanki reguły

$$
\frac{\Gamma,\neg\varphi \vdash \bot}{\Gamma \vdash \varphi}\;(\mathrm{RAA})
$$

czyli takich, których tezą jest \(\bot\) i które mają w kontekście formułę będącą negacją.

Słowniki przechowujemy w dwóch wariantach: pierwszy mapuje klucze na listy odpowiednich
indeksów, drugi z kolei na zbiory. Służy to optymalizacji różnych typów operacji.
Listy są wygodne do losowania, natomiast zbiory do szybkiego sprawdzania przynależności i wyznaczania przecięć.

Dopisanie linii dla każdej reguły wnioskowania jest w implementacji obsługiwane
osobno, ponieważ poszczególne reguły wymagają innych warunków doboru przesłanek
oraz, w niektórych przypadkach, dodatkowych argumentów. Szczegółowe omawianie
wszystkich tych przypadków nie jest jednak konieczne do zrozumienia ogólnej idei
generatora. Aby pokazać sposób wykorzystania opisanych słowników, przyjrzymy się
regule eliminacji alternatywy, która w notacji drzewiastej ma postać
\[
\frac{
    \Gamma_1 \vdash \varphi \vee \psi
    \qquad
    \Gamma_2,\varphi \vdash \rho
    \qquad
    \Gamma_3,\psi \vdash \rho
}{
    \Gamma_1,\Gamma_2,\Gamma_3 \vdash \rho
}\;(\vee E).
\]
zaś w notacji tekstowej ma postać

\[
\begin{array}{l}
\texttt{j. } \phi \lor \psi \quad \dots \\
\quad \dots \\
\texttt{k. } \rho \quad \dots \\
\quad \dots \\
\texttt{l. } \rho \quad \dots \\
\quad \dots \\
\texttt{i. } \rho \quad \texttt{\(\lor\)-elimination, j, k, l}
\end{array}
\].

Aby dopisać linię przez zastosowanie tej reguły, generator musi znaleźć trzy
linie o indeksach \texttt{j}, \texttt{k} i \texttt{l} spełniające określone warunki:
konkluzją pierwszej z nich powinna być alternatywa, druga powinna mieć w kontekście
lewą stronę tej alternatywy, trzecia powinna mieć w kontekście jej prawą stronę,
a konkluzje drugiej i trzeciej linii powinny być identyczne.

Bez omawianych słowników znalezienie takiej trójki wymagałoby bardzo wielu prób.
Liczba możliwych trójek rośnie sześciennie względem liczby linii w dużym dowodzie,
a jedynie nieliczne z nich spełniają wszystkie wymagane warunki. Użycie słowników
pozwala znacząco ograniczyć przestrzeń przeszukiwania. Procedura znajdowania
odpowiednich przesłanek wygląda następująco.

Najpierw losowana jest linia z listy linii, których tezą jest alternatywa.
Oznaczmy tę tezę jako \(\varphi \vee \psi\). Następnie tworzone są dwa zbiory:
zbiór tez linii mających w kontekście \(\varphi\) oraz zbiór tez linii mających
w kontekście \(\psi\). Jeżeli przekrój tych zbiorów jest niepusty, losowana jest
z niego formuła \(\rho\). W przeciwnym wypadku procedura wraca do wyboru pierwszej
przesłanki, czyli linii z alternatywą.

Po wybraniu formuły \(\rho\) druga przesłanka losowana jest z przekroju zbioru
linii mających w kontekście \(\varphi\) oraz zbioru linii, których tezą jest
\(\rho\). Analogicznie trzecia przesłanka losowana jest z przekroju zbioru linii
mających w kontekście \(\psi\) oraz zbioru linii, których tezą jest \(\rho\).
W ten sposób otrzymujemy trójkę linii potrzebnych do zastosowania reguły
eliminacji alternatywy.

Słowniki pomocnicze oraz sparsowana reprezentacja dużego dowodu są aktualizowane
inkrementalnie. Po dopisaniu kolejnej poprawnej linii system nie parsuje ponownie
całego dowodu ani nie buduje od zera struktur pomocniczych. Zamiast tego korzysta
z reprezentacji uzyskanej w poprzedniej iteracji i uzupełnia ją o informacje
wynikające z nowej linii: jej sparsowaną postać oraz indeksy potrzebne w odpowiednich
słownikach.

Koszt tej operacji nie musi być formalnie stały, ponieważ zależy od użytych
kontenerów i od liczby miejsc, w których należy zapisać informację o nowej linii.
Jest jednak kosztem lokalnej aktualizacji, a nie ponownego przetwarzania całego
dowodu, dlatego w praktyce rośnie znacznie wolniej niż koszt każdorazowego
odtwarzania wszystkich struktur pomocniczych.

Mimo zastosowanych usprawnień, przy dużych dowodach rzędu kilkudziesięciu tysięcy
linii czas potrzebny na dopisanie kolejnej linii nadal zauważalnie rośnie. Z tego
powodu procedura generacji uruchamiana jest wielokrotnie i równolegle, w kilku
wątkach. Wszystkie wątki startują od tych samych początkowych dużych dowodów,
jednak każdy z nich pracuje na własnej kopii dowodu, rozwijając ją niezależnie od
pozostałych. Każdde procedura generacji dla jednego dużego dowodu generuje kilkanaście tysięcy liniiz.

\section{Generowanie z dużego dowodu przykładów treningowych}

Same duże dowody nie są bezpośrednio użyteczne jako dane treningowe dla modeli
językowych. Wynika to z kilku powodów.

Po pierwsze, ich format różni się od formatu przykładów, które model ma później
generować. W przypadku rozbić celu format rozbić i dużego dowodu jest całkowicie różny, 
w przypadku dowodów z założeniami w dużych dowodach brakuje lini kodującej cel
oraz linii kodujących założenia.

Po drugie, linie dużego dowodu powstają w sposób losowy. Oznacza to, że chociaż
cały ciąg jest formalnie poprawny, to nie stanowi uporządkowanej odpowiedzi na
konkretny problem dowodowy. Nie zawiera więc bezpośrednio informacji o tym, jakich
reguł i argumentów należy używać, aby rozwiązać zadany sekwent.

Po trzecie, duże dowody są zbyt długie, aby mogły być w całości przetwarzane przez
model językowy. Ich długość znacząco przekracza długość kontekstu obsługiwaną przez
wytrenowane modele. W rezultacie model, widząc późniejszą linię dowodu, nie miałby
już dostępu do wielu wcześniejszych linii, do których ta linia się odwołuje. Uczenie
na takich przykładach nie pozwalałoby więc skutecznie nauczyć się doboru argumentów
dla poszczególnych reguł.

Z tego powodu duże dowody musiały zostać przekształcone do postaci użytecznej
dla treningu. Proces ten można podzielić na trzy zasadnicze etapy:

\begin{enumerate}[itemsep=0.4em, topsep=0.4em, parsep=0pt]
    \item \textbf{Wyłuskiwanie dowodów tautologii.}
    Z dużego dowodu wybierane są linie kodujące sekwenty o pustym kontekście.
    Każda taka linia odpowiada dowodowi pewnej tautologii. Następnie odtwarzany
    jest fragment dużego dowodu potrzebny do uzasadnienia tej linii, a otrzymane
    dowody podlegają dalszej selekcji.

    \item \textbf{Wyłuskiwanie dowodów z założeniami z dowodów bez założeń i formatowanie powstałego w ten sposób korpusu}
    Z otrzymanych dowodów tautologii dla ich poszczególnych linii generowane są dowody z założeniami oraz
    rozbicia.

    \item \textbf{Wyłuskiwanie rozbić problemów na podproblemy z dowodów bez założeń i formatowanie powstałego w ten sposób korpusu}
\end{enumerate}

\subsection{Wyłuskiwanie dowodów tautologii}

Pierwszym etapem wyłuskiwania dowodów tautologii jest sparsowanie dużego dowodu.
Następnie znajdujemy wszystkie linie, których sekwenty po sparsowaniu mają puste
konteksty (formuły zapisane w tych liniach są tautologiami). 

Następnie tworzymy zbiór linii użytych w dowodzie tej tautologii.
początkowo zbiór ten zawiera tylko linię, którą chcemy udowodnić, a następnie dodajemy do
tego zbioru wszystkie linie, do których odwołuje się ta linia. Potem analogicznie dodajemy linie,
do których odwołują się nowo dodane linie, i
powtarzamy ten proces, aż otrzymany zbiór linii będzie zamknięty ze względu na
odwołania. 

Na końcu wybrane linie sortujemy zgodnie z ich kolejnością w dużym dowodzie, 
łączymy je w jeden dowód i przenumerowujemy je tak, aby indeksy kolejnych linii 
były kolejnymi liczbami naturalnymi, zaczynając od \texttt{1}.

W tej fazie dokonujemy również wstępnej selekcji dowodów które nas interesują.

Odrzucamy:
\begin{enumerate}[itemsep=0.4em, topsep=0.4em, parsep=0pt]
    \item \textbf{Większość dowodów mających mniej niż 3 linie (zostawiamy losowo około 0.01 ich początkowej liczby).} W szczególności odrzucamy większość jednolinijkowych "dowodów" 
    tautologii przy użyciu reguł \texttt{TND} i $\top$-\texttt{introduction}. 
    \item \textbf{Dowody tautologii zawierających symbole \(\bot\) i \(\top\) (zostawiamy losowo około 0.01 ich początkowej liczby).}
        Tautologie te początkowo miały dość dużą nadreprezentację, a ich dowody często
        sprowadzały się do mechanicznego wykorzystania własności stałych logicznych, przez
        co wnosiły niewiele informacji o ogólnych schematach dowodzenia.
    \item \textbf{Wszystkie dwody mające więcej niż 1024 znaki.} Odrzucono je, ponieważ rozmiar kontekstu trenowanego modelu językowego wynosił \textbf{1024}
\end{enumerate}

W wygenerowanym zbiorze pojawiało się wiele dowodów tych samych tautologii lub
tautologii różniących się jedynie nazwami zmiennych. Aby usunąć tego typu
powtórzenia, dowody grupowano ze względu na znormalizowaną postać dowodzonej
tautologii. W postaci znormalizowanej pierwsza zmienna występująca w formule od
lewej była zastępowana przez \texttt{x1}, druga przez \texttt{x2}, trzecia przez
\texttt{x3} itd. Następnie dla każdej grupy wybierano dowód o najmniejszej liczbie
linii, a pozostałe dowody odrzucano ze zbioru z którego będziemy generować korpusy uczące.

\subsection{Tworzenie dowodów z założeniami}
Z otrzymanych dowodów tautologii generuję następnie przykłady dowodów
sekwentów z założeniami. Każda linia poprawnego dowodu koduje pewien sekwent
\(\Gamma \vdash \varphi\), a poprzedzające ją linie zawierają jego dowód.
Będziemy więc odtwarzać dowód \(\varphi\) w kontekście \(\Gamma\).

Dla każdej linii dowodu odtwarzano najpierw fragment potrzebny do jej
uzasadnienia. Procedura ta przypomina procedurę ekstrakcji dowodów tautologii
z dużego dowodu. Zaczynała ona od wybranej linii, a następnie dodawała wszystkie
linie, do których odwołuje się jej uzasadnienie. Następnie analogicznie dodawano
linie, do których odwołują się nowo dodane linie. Proces powtarzano aż do
otrzymania zbioru linii zamkniętego ze względu na odwołania. W ten sposób
z całego dowodu wycinano poddowód wystarczający do uzasadnienia wybranego
sekwentu. Po wyborze odpowiednie linie były ze sobą sklejane, a ich numeracja
oraz numeracja odwołań była zmieniana tak, aby numery kolejnych linii były
kolejnymi liczbami naturalnymi zaczynając od \texttt{1}.

Następnie tworzono zapytanie zgodne z formatem dowodów z kontekstem, używanym
w korpusie. Zapytanie rozpoczyna się linią, w której po \texttt{Prove that:}
zapisywana jest formuła, którą chcemy udowodnić. Jeżeli zbiór założeń jest
niepusty, to w kolejnej linii po tekście \texttt{assuming that:} zapisywane jest
pierwsze założenie, a kolejne założenia są zapisywane po przecinku w następnych
liniach.

Po złączeniu części odpowiadającej za definicję problemu i części zawierającej
dowód otrzymujemy dowód w formacie, który mógłby zostać użyty w korpusie.
Niestety taki dowód wciąż może zawierać liczne nadmiarowe linie, co utrudniałoby
modelowi zrozumienie, które kroki faktycznie zbliżają nas do udowodnienia
sekwentu. W tym celu stosujemy kilka uproszczeń.

Naturalnym pierwszym kierunkiem
uproszczeń było skorzystanie z jawnego kontekstu problemu. Jeżeli w dowodzie
pojawiała się linia dowodząca formuły należącej do tego kontekstu, zastępowano
jej uzasadnienie odwołaniem do kontekstu. W zależności od tego, czy linia miała
pusty czy niepusty kontekst dodatkowy, używano reguły \texttt{context} dla
argumentu \texttt{0} albo reguły \texttt{contextWeak}, również dla argumentu
\texttt{0}. Linia \texttt{0} została zdefiniowana wcześniej jako niezapisana
wprost linia, której kontekst jest równy kontekstowi bazowemu. Po tych podmianach
ponownie usuwano linie, które przestały być potrzebne, a cały fragment
przenumerowywano. W rezultacie powstawał krótszy i bardziej regularny dowód
z założeniami, lepiej nadający się do treningu modelu.

Może się również zdarzyć, że podczas dowodzenia sekwentu z założeniami jakaś
linia dowodu, która nie jest ostatnią linią, także koduje sekwent, który chcemy
udowodnić. Ponieważ dowód może zakończyć się na pierwszej linii kodującej żądany
sekwent, usuwano wszystkie linie znajdujące się po pierwszym wystąpieniu tego
sekwentu.


Kolejne uproszczenie bazuje na następującej intuicji. Zdefiniujmy indukcyjnie
zbiór linii ważnych. Niech \(n\) oznacza numer ostatniej linii dowodu \(P\), zaś
\(\operatorname{Ref}(k,P)\) oznacza zbiór linii, do których odwołuje się linia
\(k\) w dowodzie \(P\). Przyjmujemy
\[
W_{1,P} = \{n\}.
\]
Następnie definiujemy
\[
W_{i+1,P}
=
\bigcup_{k \in W_{i,P}} \operatorname{Ref}(k,P).
\]
Ostatecznie niech
\[
W_P = \bigcup_{i \in \mathbb{N}} W_{i,P}.
\]
Zbiór \(W_P\) zawiera więc ostatnią linię dowodu, linie, do których ona się
odwołuje, linie, do których odwołują się te linie itd. Innymi słowy są to linie,
które są potrzebne do uzasadnienia końcowego sekwentu.

Zauważmy, że dowód z założeniami \(P\) pozostanie poprawnym dowodem żądanego
sekwentu, jeżeli usuniemy z niego wszystkie linie nienależące do \(W_P\), a w
pozostałych liniach zmienimy numerację linii i odwołań tak, aby numery kolejnych
linii były kolejnymi liczbami naturalnymi. Linie spoza \(W_P\) nie są bowiem
wykorzystywane w uzasadnieniu ostatniej linii dowodu.

Konkluzje reguł są często w pewnym sensie podobne do ich przesłanek. Na przykład
mogą mieć takie same konteksty, konteksty różniące się jedną formułą, albo w
przypadku reguły \texttt{weakening} tę samą konkluzję. Można więc przypuszczać,
że w niektórych miejscach dowodu da się zmienić odwołanie do danej linii na
odwołanie do innej, wcześniejszej linii, nie niszcząc poprawności dowodu. W ten
sposób, zmieniając numery odwołań, możemy wpływać na rozmiar zbioru \(W_P\). Jeżeli
po takiej zmianie część późniejszych linii przestanie należeć do \(W_P\), będzie
można usunąć je z dowodu.

Znalezienie dokładnego algorytmu, który wybierałby odwołania prowadzące do
najmniejszego możliwego zbioru \(W_P\), wydaje się trudne. Dlatego zastosowano
prostą heurystykę: jeżeli w każdym odwołaniu będziemy starali się odwoływać do
jak najwcześniejszej linii, to powinno to sprzyjać zmniejszaniu liczby późnych
linii należących do \(W_P\). W szczególności późniejsze linie często są wynikiem
wielu wcześniejszych kroków dowodowych, z których każdy posiada własne odwołania.
Obecność takiej linii w \(W_P\) może więc pociągać za sobą obecność wielu innych
linii.

Aby ująć to dokładniej, dla linii \(k\) zdefiniujmy
\[
\operatorname{Ref}_1(k,P) = \operatorname{Ref}(k,P),
\]
\[
\operatorname{Ref}_{i+1}(k,P)
=
\bigcup_{j \in \operatorname{Ref}_i(k,P)} \operatorname{Ref}(j,P),
\]
oraz
\[
\overline{\operatorname{Ref}}(k,P)
=
\bigcup_{i \in \mathbb{N}} \operatorname{Ref}_{i}(k,P).
\]
Zbiór \(\overline{\operatorname{Ref}}(k,P)\) zawiera więc linie, od których linia
\(k\) zależy pośrednio i bezpośrednio. Intuicyjnie, im większy jest indeks \(k\), tym większa
jest szansa, że \(\overline{\operatorname{Ref}}(k,P)\) będzie duży. Jeżeli taka
późna linia należy do \(W_P\), to razem z nią w \(W_P\) znajdą się również linie
należące do \(\overline{\operatorname{Ref}}(k,P)\). Jest to kolejny powód, dla
którego preferowane są odwołania do linii pojawiających się w dowodzie możliwie
wcześnie.

Sama procedura zmian odwołań, ze względu na niewielkie rozmiary treningowych
dowodów, jest dość prosta. Dla każdego odwołania w dowodzie sprawdzamy, czy
zmiana wskazywanego indeksu na kolejne liczby naturalne zaczynając od \(0\) nie psuje
poprawności dowodu. Jeżeli po takiej zmianie dowód nadal przechodzi weryfikację,
podmiana jest zachowywana. Jako poprawność rozumiemy tutaj zarówno to, że dowód
pozostaje ciągiem poprawnych kroków, jak i to, że nadal dowodzi żądanego sekwentu.
Po dokonaniu wszystkich takich zmian (początkowo zmieniamy odwołania w ostatniej linii, 
następnie idziemy w górę) z dowodu usuwane są linie nienależące do \(W_P\), a
następnie zmieniana jest numeracja linii i odwołań tak, aby indeksy linii były
kolejnymi liczbami naturalnymi zaczynając od \texttt{1}. 

\subsection{Tworzenie rozbić sekwentów}


Podobnie jak w przypadku generacji dowodów z założeniami, rozbicia
generujemy z poszczególnych linii dowodu bez założeń. 

Rozbicie składa się z problemu głównego, reguły użytej do jego rozbicia oraz
listy podproblemów odpowiadających przesłankom tej reguły. Jeżeli dana linia
dowodu koduje sekwent \(P\) i została uzyskana za pomocą reguły \(R\) z linii
kodujących przesłanki \(P_1,\ldots,P_n\), to możemy odczytać z niej rozbicie
problemu \(P\) na podproblemy \(P_1,\ldots,P_n\) za pomocą reguły \(R\).

Jak zdefiniowano w sekcji~\ref{sec:reprezentacja-dowodow}, sparsowany dowód
jest słownikiem mapującym indeksy linii na odpowiadające im sparsowane linie.
Każda sparsowana linia zawiera sekwent zapisany w danej linii, regułę użytą do
jej uzasadnienia oraz indeksy linii, do których odwołuje się to uzasadnienie.
Mając więc sparsowany dowód oraz wybraną linię tego dowodu, możemy pobrać
problem główny i regułę bezpośrednio z tej linii, natomiast podproblemy
otrzymać przez odwołanie się do wcześniejszych linii wskazanych przez jej
argumenty.

Należy jednak zaznaczyć, że nie każdy indeks występujący w zapisie uzasadnienia
linii odpowiada rzeczywistej przesłance rozbicia. Jak wyjaśnialiśmy wcześniej w przypadku niektórych reguł
część indeksów odwołuje się do linii które nie kodują przesłanek ale które są konieczne, żeby 
odczytać jakiś argument reguły który nie wynika bezpośrednio z przesłanek.

Następnie rozbicia zostały przekonwertowane do formy tekstowej gdzie pierwsza linia składa się z napisu \texttt{Smash: } po
którym następuje zapis sekwentu który rozbijamy, druga linia składa się z napisu \texttt{With: } po którym następuje nazwa 
reguły używanej do rozbicia, a zaś trzecia napis \texttt{To: } po którym jeżeli lista podproblemów na które rozbijamy
sekwent jest niepusta jest zapisana pierwsza przesłanka, z kolei kolejne są zapisywane po przecinku w następnych liniach.

Wadą wygenerowanego korpusu rozbić było to, że konkluzją zdecydowanej większości rozbić przy pomocy reguły 
$$\frac{\Gamma_1,\varphi \vdash \psi \qquad \Gamma_2,\psi \vdash \varphi}
     {\Gamma_1,\Gamma_2 \vdash \varphi \leftrightarrow \psi}\;(\leftrightarrow I)$$

formuły $\varphi$ i $\psi$ były identyczne, co przy praktycznym szukaniu dowodów sekwentów logicznych zdarza się
rzadko. W tym celu dodałem do korpusu rozbicia przy pomocy reguły wyprowadzenia równoważności wygenerowane przez 
osobny algorytm.

Wartościowaniem zbioru zmiennych nazwijmy funkcję, która zmiennym z tego zbioru przypisuje wartości $\top, \bot$.
Dla zbioru $n$ zmiennych istnieje $2^n$ wartościowań.

Przyjmijmy, że mamy skończony $n$-elementowy zestaw zmiennych $A$, które będą występować w rozważanych formułach i kontekstach.

Jako $$\operatorname{Sat}(X)$$ zdefiniujmy zbiór wszystkich wartościowań zmiennych z $A$, które spełniają $X$
jeśli $X$ jest formułą i zbiór wszystkich wartościowań zmiennych z $A$, które spełniają wszystkie 
formuły w $X$ jeśli $X$ jest kontekstem.

Zauważmy, że $$\operatorname{Sat}(\Gamma,\varphi) = \operatorname{Sat}(\Gamma) \cap \operatorname{Sat}(\varphi)$$

Istotnie, jeżeli jakieś wartościowanie należy do $\operatorname{Sat}(\Gamma,\varphi)$ to musi spełniać zarówno formułę
$\varphi$ jak i jak i formuły w $\Gamma$ czyli należeć zarówno do $\operatorname{Sat}(\varphi)$ jak i do $\operatorname{Sat}(\Gamma)$, czyli należeć do przecięcia tych zbiorów.
W drugą stronę, jeżeli wartościowanie należy do $\operatorname{Sat}(\Gamma) \cap \operatorname{Sat}(\varphi)$ to spełnia zarówno formułę $\varphi$ jak i formuły 
z $Gamma$, czyli należy do $\operatorname{Sat}(\Gamma,\varphi)$.

Zauważmy również, że 
$$\Gamma \vdash \phi \leftrightarrow \operatorname{Sat}(\Gamma)\subseteq \operatorname{Sat}(\varphi) $$

Istotnie zapis $\Gamma \vdash \phi$ oznacza, że z formuł z $\Gamma$ da się dowieść $\varphi$, co w przypadku rachunku zdań w logice klasycznej jest równoważne
temu, że jeżeli dla jakiegoś prawdziwe są wszystkie formuły z $\Gamma$ to dlatego wartościowania prawdziwa jest formuła $\varphi$, co jest równoważne z tym, że 
jeżeli jakieś wartościowanie należy do $\operatorname{Sat}(\Gamma)$ musi należeć do $\operatorname{Sat}(\varphi)$, czyli, że $\operatorname{Sat}(\Gamma)$
zawiera się w $\operatorname{Sat}(\varphi)$.

Algorytm wygenerowania rozbicia przy pomocy
$$\frac{\Gamma_1,\varphi \vdash \psi \qquad \Gamma_2,\psi \vdash \varphi}
     {\Gamma_1,\Gamma_2 \vdash \varphi \leftrightarrow \psi}\;(\leftrightarrow I)$$

będzie polegał, na znalezieniu przesłanek tego rozbicia, by później wywnioskować z nich konkluzję, a potem złączyć konkluzję, przesłanki i rozbicie w sparsowane rozbicie,
które później sprowadzimy do postaci tekstowej.

Żeby znaleźć przesłanki reguły musimy znaleźć konteksty $\Gamma_1$ i $\Gamma_2$ oraz formuły $\varphi$, $\psi$
takie, że 
$$\Gamma_1,\varphi \vdash \psi$$
oraz
$$\Gamma_2,\psi \vdash \varphi$$
Ze wcześniejszych obserwacji wynika, że warunki te są równoważne z
$$\operatorname{Sat}(\Gamma_1) \cap \operatorname{Sat}(\varphi) \subseteq \operatorname{Sat}(\psi)$$
oraz
$$\operatorname{Sat}(\Gamma_2) \cap \operatorname{Sat}(\psi) \subseteq \operatorname{Sat}(\varphi)$$

Na początku nasz algorytm losuje konteksty $\Gamma_1$ i $\Gamma_2$, a następnie w pętli losuje parę formuł $\varphi$ 
i $\psi$ tak aby spełniały powyższe warunki. Taka para zawsze istnieje dla dowolnych kontekstów 
$\Gamma_1$ i $\Gamma_2$ np zawsze warunki te spełnia para tautologii, lub formuł sprzecznych. Samo losowanie 
formuł wydaje się operacją, która wymaga bardzo wielu powtórzeń, jednak ograniczając liczbę różnych zmiennych które 
mogą występować łącznie w generowanych omawianym algorytmem kontekstach i formułach do $6$ udaje się wygenerować przy pomocy opisywanej metody wiele rozbić w stosunkowo krótkim czasie.

Kolejnym problemem jest to, że nie zawsze wymagane właściwości przesłanek i konkluzji w wygenerowanych na tym etapie danych
przekładają się na proste tekstowe wzorce w łatwo identyfikowane przez modele językowe. Na przykład weźmy regułę

$$
\displaystyle
\frac{\Gamma_1 \vdash \varphi \vee \psi \qquad \Gamma_2,\varphi \vdash \rho \qquad \Gamma_3,\psi \vdash \rho}
     {\Gamma_1,\Gamma_2,\Gamma_3 \vdash \rho}\;(\vee E)
$$

Przykład jej użycia w notacji drzewiastej mógłby wyglądać następująco

$$
\displaystyle
\frac{a \wedge h \vdash a \vee b \qquad a , c \wedge d, f \vdash c \qquad c,b,e \vdash c}
     {c \wedge d, e, c, f, a\wedge h \vdash c}\;(\vee E)
$$
co w notacji zapisu rozbić przekładało by się na następujący tekst:
\[
\begin{array}{l}
\texttt{Smash: c}\wedge \texttt{d, e, c, f, a } \vdash \texttt{ c}\\
\texttt{With: DisjunctionElimination} \\
\texttt{To: a}\wedge\texttt{h } \vdash \texttt{ a} \vee \texttt{b,}\\
\texttt{a, c}\wedge\texttt{d, f } \vdash  \texttt{ c,} \\
\texttt{c, b, e }\vdash \texttt{ c}
\end{array}
\]

Z podanego napisu trudno wywnioskować na pierwszy rzut oka jakikolwiek regularność. Jednocześnie zauważmy, że przy generacji rozbicia przy pomocy
eliminacji alternatywy, generujemy również poprawne rozbicia jeśli założymy, że $\Gamma_1$, $\Gamma_2$, $\Gamma_3$ są takie same i równe ich sumie 
a wtedy jeżeli zachowamy odpowiednią kolejność wypisywania przesłanek sekwentów (co zapewniamy odpowiednią implementacją działań na kontekstach) możemy zamienić powyższy zapis rozbicia na 
\[
\begin{array}{l}
\texttt{Smash: c}\wedge \texttt{d, e, c, f, a}\wedge\texttt{h } \vdash \texttt{ c}\\
\texttt{With: DisjunctionElimination} \\
\texttt{To: c}\wedge \texttt{d, e, c, f, a}\wedge\texttt{h }\vdash\texttt{ a}\vee\texttt{b}\\
\texttt{c}\wedge \texttt{d, e, c, f, a}\wedge\texttt{h, a} \vdash  \texttt{ c,} \\
\texttt{c}\wedge \texttt{d, e, c, f, a}\wedge\texttt{h, b}\vdash \texttt{ c}
\end{array}
\]
gdzie wzorzec jest dużo bardziej zauważalny. Konteksty konkluzji i przesłanek są niemal identyczne i w przypadku dwóch ostatnich przesłanek różnią się jedynie jedną formułą na końcu.

Aby możliwie uprościć modelowi znajdowanie porządanych wzroców nalogiczne przekształcenia, których celem jest możliwie regularny zapis rozbicia definiujemy dla rozbić przy pomocy pozostałych reguł. 
Transformacje te polegają, na dodaniu do kontekstów przesłanek rozbicia pewne,
być może puste, zbiory formuł i ustawieniu formuł w kontekstach w odpowiedniej kolejności. Transformacja to odbywa się w taki sposób, aby
jeżeli początkowe rozbicie było poprawne (to znaczy pasowało do schematu reguły i miało w przesłankach sekwenty tautologiczne) otrzymane po transformacji
rozbicie nadal było poprawne a także miało tą samą regułę i konkluzję co przed transformacją. Transformacje te stosujemy do wszystkich rozbić w korpusie.


Ostatnim etapem formatowania korpusu rozbić jest modyfikacja częstości występowania rozbić tworzonych przez poszczególne reguły. W korpusie który wygenerowaliśmy do tej pory
częstość występowania rozbić generowanych przez niektóre reguły znacznie przewyższała częstości innych, przez co model często używał częściej występujących reguł bez
nawet w sytuacjach w których nie były one użyteczne. Co więcej ideą rozbić jest to, by rozbijały problemy na podproblemy prostsze, a dobrą heurystyką na to czy dany problem jest prostszy
czy trudniejszy jest długość jego konkluzji. Niektóre reguły rozbijają swoją konkluzję na przesłanki których konkluzje są dłuższe 
niż konkluzja konkluzji reguły (chcielibyśmy mieć ich mniej w korpusie) na przykład
$$\frac{\Gamma \vdash \varphi}{\Gamma \vdash \varphi \vee \psi}\;(\vee I_1)$$
inne zaś rozbijają konkluzję reguły na przesłanki o konkluzjach krótszych niż konkluzje konkluzji reguły (chcielibyśmy mieć ich więcej w korpusie)
na przykład
$$\frac{\Gamma_1 \vdash \varphi \qquad \Gamma_2 \vdash \psi}{\Gamma_1,\Gamma_2 \vdash \varphi \wedge \psi}\;(\wedge I)$$
 
W tym celu rozbicia w korpusie grupujemy ze względu na reguły w nich występujące, a następnie wielokrotnie losujemy z intuicyjnie 
ustalonym rozkładem regułę generującą rozbicie, a następnie jednostajnie losujemy rozbicie ze zbioru rozbić generowanych przez
tą regułę i dodajemy je do nowego korpusu. W ten sposób powstaje ostateczna wersja danych treningowych dla modelu odpowiadającego za rozbicia.

\section{Losowe rozszerzanie dużego dowodu}

\section{Wyłuskiwanie dowodów tautologii}

\section{Normalizacja formuł i usuwanie powtórzeń}

\section{Generowanie problemów pomocniczych}

\section{Generowanie smashy}

\section{Normalizacja kontekstów w smashach}

\section{Równoważność jako przypadek specjalny}

\section{Walidacja i czyszczenie korpusu}

\section{Podsumowanie}

\chapter{Trenowanie i kożystanie z modelu (NanoGPTBundle i te wszystkie funkcje z nim związane !!!!!!!!!!!!!!!!!!!!!!!!!!!!!!!!, podkreślam WSZYSTKIE(liczenie prawdopodobieństwa, dodawanie tokenu, linii, trenowanie, dotrenowanie, ładowanie)!!!!)}


%%%%% BIBLIOGRAFIA


%\begin{thebibliography}{1}
%\bibitem{example} \ldots
%\end{thebibliography}

\end{document}
